\documentclass[UTF8, scheme=chinese, oneside]{ctexbook}

\usepackage{pdflscape}% 引入横向排版宏包

% 让正文区域占整个页面纸张面积的 80%
\usepackage[scale=0.8]{geometry}

\usepackage{xeCJK}
\xeCJKsetup{
  AutoFallBack = true,
  AutoFakeBold = true,
  PunctStyle = kaiming,
  CheckFullRight = true,
  CheckSingle = false,
  AllowBreakBetweenPuncts = true,
  CJKmath = true,
}

% 完整的中文章节格式设置
\ctexset{
    % 章设置
    % chapter = {
    %     format = \huge\bfseries\centering,
    %     nameformat = {},
    %     titleformat = {},
    %     name = {第,章},
    %     number = \chinese{chapter},
    %     aftername = \quad,
    %     beforeskip = 20pt,
    %     afterskip = 20pt,
    % },
    % 节设置
    section = {
        format = \Large\bfseries,
        nameformat = {},
        titleformat = {},
        name = {第,节},
        number = \chinese{section},
        aftername = \quad,
    },
    % 小节设置
    subsection = {
        format = \large\bfseries,
        nameformat = {},
        titleformat = {},
        name = {第,句},
        number = \chinese{subsection},
        aftername = \quad,
    },
    % 子小节设置
    subsubsection = {
        format = \normalsize\bfseries,
        nameformat = {},
        titleformat = {},
        name = {第,子小节},
        number = \chinese{subsubsection},
        aftername = \quad,
    }
}

\xeCJKDeclareSubCJKBlock{Ext-A}{ "3400 -> "4DB5 }
\xeCJKDeclareSubCJKBlock{Ext-B}{ "20000 -> "2A6DF }
\xeCJKDeclareSubCJKBlock{CircNum}{ "2460 -> "2473 }
\xeCJKDeclareSubCJKBlock{Symbols}{ "2600 -> "27FF }


\setCJKmainfont[%
BoldFont={FZYouHei_GBK 513B},   %
ItalicFont={FZNewKai-Z03},
% Ext-A={仓耳今楷05-W04},
Ext-B={SimSun-ExtB},
CircNum={IosevkaTerm NFM},
Symbols={IosevkaTerm NFM}]
{NSimSun}
\newCJKfontfamily[li]\hmls{FZZJ-HMLSJF}
\newCJKfontfamily[kai]\fzxk{FZNewKai-Z03}
\newCJKfontfamily[song]\ssb{SimSun-ExtB}

\usepackage{xcolor}
\usepackage{pifont}

% 引入 tikz 及相关路径变形库（用于画波浪线）
\usepackage{tikz}
\usetikzlibrary{decorations.pathmorphing}

% 自定义双波浪线命令 \uuwave
\newcommand{\uuwave}[1]{%
  \tikz[baseline=(char.base)]{
    % 将文字放入一个无边距的节点中
    \node[inner sep=0pt, outer sep=0pt] (char) {#1};
    % 画第一条波浪线 (距离文字稍近)
    \draw[decorate, decoration={snake, amplitude=0.3mm, segment length=2mm}]
      ([yshift=-1.5pt]char.south west) -- ([yshift=-1.5pt]char.south east);
    % 画第二条波浪线 (在第一条下方)
    \draw[decorate, decoration={snake, amplitude=0.3mm, segment length=2mm}]
      ([yshift=-3.0pt]char.south west) -- ([yshift=-3.0pt]char.south east);
  }%
}

% 引入下划线宏包，normalem 参数防止破坏原有的 \emph{} 斜体命令
\usepackage[normalem]{ulem}

% 引入颜色宏包（可选，用于让中文标注颜色变浅一点）
\usepackage{xpinyin}

% 设置段落缩进和行距
\setlength{\parindent}{0pt}
\linespread{1.5}

% ==========================================
% 自定义语法标注命令 \markword
% 用法: \markword[线型命令]{英文内容}{中文标注}
% 默认线型为单实线 \uline。如果不想要线，可以传入 \mbox
% ==========================================
\newcommand{\markword}[3][\uline]{%
  \begin{tabular}[t]{@{}c@{}}%
    #1{#2} \\[-0.2ex]% 调整英文和中文之间的垂直间距
    \small\color{darkgray} #3% 中文标注字号设小，颜色设为深灰
  \end{tabular}%
}

% 下面两个数学符号的package是互相冲突的！
\usepackage{amsmath}
\usepackage{mathabx}
% \usepackage{fdsymbol}

\usepackage{changepage}

% 表格
\usepackage{enumitem}
\usepackage{tabularray}
\UseTblrLibrary{varwidth}
% 1. 修改当前页底部的提示文字 (默认是 Continued on next page)
\DefTblrTemplate{contfoot-text}{default}{下页继续}

% 2. 修改下一页顶部的提示文字 (默认是 (Continued))
\DefTblrTemplate{conthead-text}{default}{(续表)}

\title{\fzxk {\huge 雅思语法学习笔记}}

\begin{document}
\maketitle

\tableofcontents
\newpage

% %% -*- coding: utf-8 -*-

\chapter{语法成分}
\newpage

% 使用 landscape 环境包裹你的表格
\begin{landscape}
  \pagestyle{empty}
  \begin{center}
    \begin{longtblr}[
      label = none,
      ]{
      vlines, hlines,
      % 移除全局的 halign=c，改为在每一列单独设置，这样更灵活
      columns={valign=m},
      column{1}={0.15\linewidth, halign=c, bg=azure9},
      column{2}={0.35\linewidth, halign=l},
      % 关键修改：给第三列指定宽度（X代表填满剩余所有空间），并设为左对齐
      column{3}={0.5\linewidth, halign=l},
      rows={abovesep=5pt},
      row{1}={font=\bfseries, bg=azure7, halign=c}, % 表头居中
      % 关键修改：处理单元格内的 list 环境
      measure=vbox,
      stretch=-1,
      }
      语法成分                                             & 定义                                                                                                                    & 注释及示例                                                                                                                                                                                \\
      \textbf{宾语} (The Object) —— “动作的承受者”                  & 宾语接在及物动词（Transitive Verb）之后，表示动作涉及的对象或受动者。                                                   & $主语 (A) \xrightarrow{动作} 宾语 (B)$. 注意：A 对 B 做了某事，A $\neq$ B                                                                                                                    \\
      \textbf{表语} (The Predictive) —— “主语的等号”                & 表语接在系动词（Linking Verb）之后，用来描述主语的身份、特征、状态或性质。                                              & $主语 (A) = 表语 (A')$. 表语其实就是主语的另一种描述，A $=$ A'.                                                                                                                           \\
      \textbf{系动词} (Linking Verbs)                               & 系动词本身没有具体的动作含义，它的作用像一把“等号” (=)，用来连接主语和描述主语状态的信息。                              & \begin{description}[nosep]
                                                                                                                                       \item[最常见的系动词] be (am, is, are, was, were)
                                                                                                                                       \item[感官系动词] look (看起来), sound (听起来), smell (闻起来), feel (摸起来/感觉)
                                                                                                                                       \item[变化/状态系动词] become (变成), stay (保持), remain (依然是)
                                                                                                                                     \end{description}                                                                                                                                                                                                                                           \\
      \textbf{后置定语} (Post-modifier / Post-positional Attribute) & 它像胶水一样紧紧粘在名词后面，用来限定或修饰这个名词“是哪一个”或“什么样的”。它属于名词短语（Noun Phrase）内部的一部分。 & The engineer [working at this company] is my friend. \newline \textit{working at this company} 是后置定语。它定义了是“哪位”工程师。如果删掉它，句子的主语依然完整，但范围变模糊了。                \\
      \textbf{主语补足语} (Subject Complement)                      & 它是用来补充说明主语的状态或身份的。它不是名词的“零件”，而是通过系动词（Linking Verb）或某些特定动词与主语建立联系。    & 它属于句子骨干（Clause Structure）的一部分。删掉后，句子意思不完整（如果是系动词后）。 \newline 如：The chip remains \textit{[untested]}. \newline \textit{untested} 是主语补足语。它说明了芯片目前的状态。 \\
      \textbf{同位语}（Appositive）                                 & 两个名词（或名词性短语）并列放在一起，支撑或解释同一个对象。第二个名词是对第一个名词的身份重塑或进一步补充。            & $[Noun A], [Noun B], ...$  （其中 $A = B$）. 注意，同位语通常用来解释名词的内容 (What it is)，而定语从句通常用来描述名词的特征 (Which one)。                                              \\
    \end{longtblr}
  \end{center}
\end{landscape}

%%% Local Variables:
%%% TeX-engine: xetex
%%% TeX-master: "../IELTS_300_Sentences"
%%% End:

% LocalWords:  vlines hlines halign valign xetex abovesep nosep

% %% -*- coding: utf-8 -*-

\chapter{词性}

\section{名词和冠词}

\subsection{名词类别以及复合名词}

\begin{center}
  \begin{talltblr}[
    label = none, % Removes the "Table X:" text
    ]{
    vlines, hlines,
    % 加入支持列表环境的测量参数 (记得在导言区加载 \UseTblrLibrary{varwidth})
    % measure=vbox, stretch=-1,
    rows={halign=c, valign=m},
    columns={halign=c, valign=m},
    column{1}={0.25\linewidth, halign=c, bg=azure9},
    column{2}={0.6\linewidth, halign=l},
    % column{3}={0.35\linewidth, halign=l},
    % column{4}={0.35\linewidth, halign=l},
    row{1}={font=\large\bfseries, bg=azure9, halign=c}
    }
    形式                     & 举例                                                                                                         \\
    两个或多个词连写         & handbag(\(n\). + \(n\).), greenhouse(\(adj\). + \(n\).), output(\(adv\). + \(v\).), underground(\(prep\). + \(n\).), handwriting(\(n\). + \(v\).ing) \\
    两个或多个词用连字符连接 & good-looking(\(adj\).+\(v\).ing), fire-raiser(\(n\).+\(n\).), life-long(\(n\).+\(adj\).), summing-up(\(v\).ing+\(adv\).), mother-in-law(\(n\).+\(prep\).+\(n\).)\\
    两个或多个词分开写       & river bank(\(n\).+\(n\).), driving license(\(v\).ing+\(n\).)\\
  \end{talltblr}
\end{center}

\markword[\uuline]{\textbf{Britain}}{主语 1}
\markword{had produced}{谓语}
\markword[\uwave]{textiles}{宾语}
\markword[\mbox]{(like wool, linen and cotton),}{后置定语 (修饰 textiles)}
\markword[\mbox]{[for hundreds of years],}{时间状语}

\vspace{0.5cm} % 手动换行并增加行距

\markword[\mbox]{but}{连词}
\markword[\mbox]{[prior to the Industrial Revolution],}{时间状语}
\markword[\uuline]{\textbf{the British textile business}}{主语 2}
\markword{was}{系动词}

\vspace{0.5cm}

\markword[\uuwave]{a true ``cottage industry'',}{表语}
\markword[\mbox]{[with the work performed in small workshops or even homes by}{伴随状语}

\vspace{0.5cm}

individual spinners, weavers and dyers].

\subsection{可数名词变复数（规则变化）}

名词除了根据意义划分，还可以根据其可数性分为“可数名词”和“不可数名词”。其中可数名词在变复数时，一般情况下在词尾加-s或-es，如句中的tardigrades、temperatures和degrees都是可数名词的复数形式，直接在词尾加-s。
\begin{center}
  \begin{talltblr}[
    label = none, % Removes the "Table X:" text
    ]{
    vlines, hlines,
    % 加入支持列表环境的测量参数 (记得在导言区加载 \UseTblrLibrary{varwidth})
    % measure=vbox, stretch=-1,
    rows={halign=c, valign=m},
    columns={halign=c, valign=m},
    column{1}={0.1\linewidth, halign=c, bg=azure9},
    column{2}={0.35\linewidth, halign=l},
    column{3}={0.4\linewidth, halign=l},
    % column{4}={0.35\linewidth, halign=l},
    row{1}={font=\large\bfseries, bg=azure9, halign=c}
    }
    类别&构成方法&举例\\
    \SetCell[r=3]{c}结尾加-s&词尾直接加-s&restaurant \(\rightarrow\) restaurants; professor \(\rightarrow\) professors\\
        &以“元音+y”结尾的名词，加-s&attorney \(\rightarrow\) attorneys; alloy \(\rightarrow\) alloys\\
        &以“元音+o”结尾的名词，加-s&studio \(\rightarrow\) studios; kangaroo \(\rightarrow\) kangaroos\\
    \SetCell[r=4]{c}结尾加-es&以 [s]、[z]、[ʃ]、[ʒ]、[ʧ]、[ʤ]等音节结尾的名词，加-es&buzz \(\rightarrow\) buzzes; bunch \(\rightarrow\) bunches; brush \(\rightarrow\) brushes; watch \(\rightarrow\) watches\\
        &以“辅音+o”结尾的名词，加 -es&volcano \(\rightarrow\) volcanoes; potato \(\rightarrow\) potatoes\\
        &词尾为-f或-fe，变f、fe 为v，加-es&shelf \(\rightarrow\) shelves; thief \(\rightarrow\) thieves\\
        &以“辅音+y”结尾的名词，变y为i加-es&strategy \(\rightarrow\) strategies; agency \(\rightarrow\) agencies\\
  \end{talltblr}
\end{center}


\markword[\mbox]{There's been}{there be 结构}
\markword[\uuline]{\textbf{some interesting research}}{主语}
\markword[\mbox]{{\large\textcolor{red}{\(\lgroup\)}}}{\textcolor{red}{定从(research)开始}}
\markword[\mbox]{which}{引导词/主语}

\vspace{0.5cm}

\markword{has found}{谓语}
\markword[\mbox]{\textcolor{blue}{\(\lgroup\)}}{\textcolor{blue}{宾从开始}}
\markword[\mbox]{that}{引导词}
\markword[\uuline]{tardigrades}{主语 1}
\markword{are capable of}{谓语 1}
\markword[\uwave]{surviving radiation and very high pressure,}{宾语 1}

\vspace{0.5cm}

\markword[\mbox]{and}{连词}
\markword[\uuline]{they}{主语 2}
\markword{are also able to withstand}{谓语 2}
\markword[\uwave]{temperatures}{宾语 2}
\markword[\mbox]{(as cold as -200 degrees centigrade),}{后置定语（temperatures）}

\vspace{0.5cm}

\markword[\mbox]{or}{连词}
\markword[\uwave]{highs}{宾语 3}
\markword[\mbox]{(of more than 148 degrees centigrade),}{后置定语（highs）}
\markword[\mbox]{(which is incredibly hot).}{定语从句}
\markword[\mbox]{\textcolor{blue}{\(\rgroup\)}}{\textcolor{blue}{宾从结束}}
\markword[\mbox]{\large{\textcolor{red}{\(\rgroup\)}}}{\textcolor{red}{定从结束}}

\subsection{可数名词变复数（不规则变化）}

可数名词变复数并非都是规则变化，还有一些特殊变化，比如句中的phenomena就是复数形式，其原形为phenomenon。
\begin{center}
  \begin{talltblr}[
    label = none, % Removes the "Table X:" text
    ]{
    vlines, hlines,
    % 加入支持列表环境的测量参数 (记得在导言区加载 \UseTblrLibrary{varwidth})
    % measure=vbox, stretch=-1,
    rows={halign=c, valign=m},
    columns={halign=c, valign=m},
    column{1}={0.25\linewidth, halign=c, bg=azure9},
    column{2}={0.55\linewidth, halign=l},
    % column{3}={0.35\linewidth, halign=l},
    % column{4}={0.35\linewidth, halign=l},
    row{1}={font=\large\bfseries, bg=azure9, halign=c}
    }
    构成方法&举例\\
    改变中间元音&goose \(\rightarrow\) geese; policeman \(\rightarrow\) policemen; analysis \(\rightarrow\) analyses; tooth \(\rightarrow\) teeth\\
    单复数同形&sheep \(\rightarrow\) sheep; salmon \(\rightarrow\) salmon\\
    结尾加-en或-ren&ox \(\rightarrow\) oxen; child \(\rightarrow\) children\\
    其他不规则形式&phenomenon \(\rightarrow\) phenomena; criterion \(\rightarrow\) criteria\\
  \end{talltblr}
\end{center}


\markword[\mbox]{[In the 1960s],}{时间状语}
\markword[\uuline]{the astronomer Gerald Hawkins}{主语}
\markword{suggested}{谓语}
\markword[\mbox]{\large\textcolor{blue}{\(\lgroup\)}}{宾从开始}
\markword[\mbox]{that}{引导词}

\vspace{0.5cm}

\markword[\uuline]{the cluster of megalithic stones}{主语}
\markword[\mbox]{operated}{谓语}
\markword{[as a form of calendar],}{方式状语}

\vspace{0.5cm}

\markword[\mbox]{[with different points corresponding to astronomical phenomenal]}{伴随状语}

\vspace{0.5cm}

\markword[\mbox]{<such as solstices, equinoxes and eclipses>}{同位语（修饰 phenomenal）}
\markword[\mbox]{(occurring at different times of the year).}{定语（修饰同位语）}
\markword[\mbox]{\large\textcolor{blue}{\(\rgroup\)}}{宾从结束}

\subsection{不可数名词}

不可数名词是指无法数清具体数量的名词，常见的不可数名词分为以下四类。
\begin{center}
  \begin{talltblr}[
    label = none, % Removes the "Table X:" text
    ]{
    vlines, hlines,
    % 加入支持列表环境的测量参数 (记得在导言区加载 \UseTblrLibrary{varwidth})
    % measure=vbox, stretch=-1,
    rows={halign=c, valign=m},
    columns={halign=c, valign=m},
    column{1}={0.15\linewidth, halign=c, bg=azure9},
    column{2}={0.35\linewidth, halign=l},
    column{3}={0.35\linewidth, halign=l},
    % column{4}={0.35\linewidth, halign=l},
    row{1}={font=\large\bfseries, bg=azure9, halign=c}
    }
    不可数类别&举例&不可数情况说明\\
    分割不变&wine, butter, honey&切开这类物质后，性质没有变化\\
    难以数清&dust, flour, pepper, hair&数量很多，很难数清楚\\
    总称概念&data, furniture, clothing&表示总称的概念，不是一个具体的事物\\
    抽象名词&health, happiness, patience, ambition, knowledge, courage&抽象的思想、观点等，无法计数\\
  \end{talltblr}
\end{center}
注意:还有一些不可数名词看上去像是可数名词的复数形式，但实际上却是不可数名词，如本句中的economics，类似的词还有news 以及 physics。

\markword[\uuline]{The book}{主语}
\markword{combines}{谓语}
\markword[\uwave]{geology, history, economics, and a lot of data}{宾语}

\vspace{0.5cm}

\markword[\mbox]{\Large\textcolor{orange}{\( [ \)}}{\textcolor{orange}{目的壮语开始}}
\markword[\mbox]{to explain}{}
\markword[\mbox]{why}{引导词}
\markword[\uuline]{business clusters}{主语}
\markword{developed}{谓语}
\markword[\mbox]{[where they did]}{地点状语从句}

\vspace{0.5cm}

\markword[\mbox]{how}{引导词}
\markword[\uuline]{the early decisions of workers and firms}{主语}
\markword{shaped}{谓语}
\markword[\uwave]{the skyline}{宾语}
\markword[\mbox]{(we see today).}{定语从句}
\markword[\mbox]{\Large\textcolor{orange}{\( ] \)}}{\textcolor{orange}{目的壮语结束}}

\subsection{名词修饰语}

想要描述不同类型名词的数量时需要注意，有些修饰语只能修饰可数名词，有些只能修饰不可数名词，有些二者均可修饰。如a picce of在本句中就用来修饰不可数名词art，表示“一件艺术品”。常用的名词修饰语还有:
\begin{center}
  \begin{talltblr}[
    label = none, % Removes the "Table X:" text
    ]{
    vlines, hlines,
    % 加入支持列表环境的测量参数 (记得在导言区加载 \UseTblrLibrary{varwidth})
    % measure=vbox, stretch=-1,
    rows={halign=c, valign=m},
    columns={halign=c, valign=m},
    column{1}={0.2\linewidth, halign=c, bg=azure9},
    column{2}={0.65\linewidth, halign=l},
    % column{3}={0.35\linewidth, halign=l},
    % column{4}={0.35\linewidth, halign=l},
    row{1}={font=\large\bfseries, bg=azure9, halign=c}
    }
    类别&举例\\
    只修饰可数名词&few (很少，几乎没有), several (几个), many (许多), a number of (大量的), scant (为数不多的，不足的), numerous (众多的，许多的), a handful of (少数，几个), a few (几个)\\
    只修饰不可数名词&little (很少，几乎没有), a little (一些，少量), a bit of (有点儿), much (许多), a large amount of (大量的), a great deal of (很多), a great quantity of (大量的), an abundance of (丰富的，充裕的)\\
    可数名词、不可数名词均可修饰&some (一些), plenty of (许多), most (大多数的), a wide range of (广泛的), substantial (大量的), a lot of/lots of (许多), enough (足够的), various (各种各样的), ample (充足的), considerable (相当多的，相当大的)\\
  \end{talltblr}
\end{center}


\markword[\mbox]{[Likewise],}{状语}
\markword[\mbox]{\Large\textcolor{orange}{\( [ \)}}{\textcolor{orange}{条件壮语开始}}
\markword[\mbox]{if}{引导词}
\markword[\uuline]{the tree}{主语}
\markword{suffers from}{谓语}
\markword[\uwave]{black knot disease}{宾语}
\markword[\mbox]{{\Large\textcolor{orange}{\( ] \)}},}{\textcolor{orange}{条件壮语结束}}

\vspace{0.5cm}

\markword[\uuline]{its value for timber}{主语 1}
\markword{decreases,}{谓语}
\markword[\mbox]{but}{连词}
\markword[\mbox]{[to a woodworker]}{对象状语}
\markword[\mbox]{(interested in making bowls),}{后置定语（woodworker）}

\vspace{0.5cm}

\markword[\uuline]{it}{主语 2}
\markword{brings}{谓语 2}
\markword[\uwave]{an opportunity}{宾语}
\markword[\mbox]{(for a unique and beautiful piece of art).}{后置定语（opportunity）}

\subsection{即可数又不可数名词}

有些名词既可数又不可数，含义会随语境而变，比如:
\begin{center}
  \begin{talltblr}[
    label = none, % Removes the "Table X:" text
    ]{
    vlines, hlines,
    % 加入支持列表环境的测量参数 (记得在导言区加载 \UseTblrLibrary{varwidth})
    % measure=vbox, stretch=-1,
    rows={halign=c, valign=m},
    columns={halign=c, valign=m},
    column{1}={0.15\linewidth, halign=c, bg=azure9},
    column{2}={0.35\linewidth, halign=l},
    column{3}={0.35\linewidth, halign=l},
    % column{4}={0.35\linewidth, halign=l},
    row{1}={font=\large\bfseries, bg=azure9, halign=c}
    }
    名词&不可数&可数\\
    paper&{含义:纸\\I need some paper to write on. 我需要一些纸来写字。}&{含义:论文;试卷\\These three papers on artificial intelligence are really insightful. 这三篇关于人工智能的论文很有见地。}\\
    room&{含义:空间\\There is no room for you to put another chair here. 这里没有空间让你再放一把椅子了。}&{含义:房间\\There are five rooms in this house.这所房子有五个房间。}\\
    light&{含义:光;光线\\The light in the room is too dim. 房间里的光线太暗了。}&{含义:灯;光源\\There are several lights in the street. 街上有几盏灯。}\\
  \end{talltblr}
\end{center}


\markword[\uuline]{The burial chamber of the tomb,}{主语}
\markword[\mbox]{{\Large\textcolor{red}{\( \lgroup \)}}}{\textcolor{red}{定从1(tomb)开始}}
\markword[\mbox]{where}{引导词}
\markword[\uuline]{the king's body}{主语}
\markword{was laid}{谓语}

\vspace{0.5cm}

\markword[\mbox]{[to rest],}{目的状语}
\markword[\mbox]{{\Large\textcolor{red}{\( \rgroup \)}}}{\textcolor{red}{定从1结束}}
\markword{was dug}{谓语}
\markword[\mbox]{[beneath the base of the pyramid],}{地点状语}

\vspace{0.5cm}

\markword[\mbox]{(surrounded by a vast maze of long tunnels)}{后置定语（chamber）}
\markword[\mbox]{{\Large\textcolor{red}{\( \lgroup \)}}}{\textcolor{red}{定从2(tunnels)开始}}
\markword[\mbox]{that}{引导词/主语}


\vspace{0.5cm}

\markword{had}{谓语}
\markword[\uwave]{rooms}{宾语}
\markword[\mbox]{[off them]}{状语}
\markword[\mbox]{[to discourage robbers].}{目的状语}
\markword[\mbox]{{\Large\textcolor{red}{\( \rgroup \)}}}{\textcolor{red}{定从2结束}}

\subsection{名词所有格}

在表达所属关系时，可以借助名词所有格。名词所有格主要分为's所有格和of所有格，有时两者会一同使用，即“双重所有格”。
\begin{center}
  \begin{talltblr}[
    label = none, % Removes the "Table X:" text
    ]{
    vlines, hlines,
    % 加入支持列表环境的测量参数 (记得在导言区加载 \UseTblrLibrary{varwidth})
    % measure=vbox, stretch=-1,
    rows={halign=c, valign=m},
    columns={halign=c, valign=m},
    column{1}={0.1\linewidth, halign=c, bg=azure9},
    column{2}={0.1\linewidth, halign=c, bg=azure9},
    column{3}={0.2\linewidth, halign=l},
    column{4}={0.45\linewidth, halign=l},
    row{1}={font=\large\bfseries, bg=azure9, halign=c}
    }
    类型&词尾&构成&举例\\
    \SetCell[r=2]{c}'s所有格&一般性词尾&词尾+'s&the scientist's research findings (科学家的研究成果)\\
        &-s, -es&词尾+'&two hours’ drive from the city center (离市中心两小时车程的地方)\\
    \SetCell[c=2]{c}of所有格&&名词+of+名词&the complexity of the economic system (经济体系的复杂性)\\
    \SetCell[r=2,c=2]{c}双重所有格&&of+名词-'s&a painting of Picasso’s (毕加索的一幅画作)\\
        &&of+名词性物主代词&a novel of his (他的一本小说)\\
  \end{talltblr}
\end{center}


\markword[\mbox]{[By the mid-1920s],}{时间状语}
\markword[\uuline]{O'keeffe}{主语}
\markword{was recognized}{谓语}

\vspace{0.5cm}

\markword[\mbox]{<as one of America's most important and successful artists>,}{主语补足语}

\vspace{0.5cm}

\markword[\mbox]{(widely known for the architectural pictures)}{后置定语（主语）}
\markword[\mbox]{{\Large\textcolor{red}{\( \lgroup \)}}}{\textcolor{red}{定从(pictures)开始}}

\vspace{0.5cm}

\markword[\mbox]{that}{引导词/主语}
\markword[\mbox]{[dramatically]}{状语}
\markword{depict}{谓语}
\markword[\uwave]{the soaring skyscrapers of New York.}{宾语}
\markword[\mbox]{{\Large\textcolor{red}{\( \rgroup \)}}}{\textcolor{red}{定从结束}}

\subsection{不定冠词a和an}

不定冠词有a和 an,a用在辅音音素前面（注意不是辅音字母），如a challenge“一项挑战”；an用在元音音素前面（注意不是元音字母），如本句中的an era“一个时代”。
\begin{center}
  \begin{talltblr}[
    label = none, % Removes the "Table X:" text
    ]{
    vlines, hlines,
    % 加入支持列表环境的测量参数 (记得在导言区加载 \UseTblrLibrary{varwidth})
    % measure=vbox, stretch=-1,
    rows={halign=c, valign=m},
    columns={halign=c, valign=m},
    column{1}={0.1\linewidth, halign=c, bg=azure9},
    column{2}={0.23\linewidth, halign=l},
    column{3}={0.55\linewidth, halign=l},
    % column{4}={0.35\linewidth, halign=l},
    row{1}={font=\large\bfseries, bg=azure9, halign=c}
    }
    用法&含义&举例\\
    \SetCell[r=5]{c}一般用法&泛指某一类人、事或物&A dog is a loyal animal. 狗是一种忠诚的动物。\\
        &首次提到某人或某物时&I saw a man walking down the street. 我看到一个男人在街上走。\\
        &表示数量“一”&Can I have a cup of coffee, please? 请给我一杯咖啡好吗?\\
        &表示“每一”&The car is running at a speed of 80 kilometers an hour. 汽车以每小时 80千米的速度行驶。\\
        &表示“又一，再一”&Can you give me a second chance? 你能再给我一次机会吗?\\
    特殊用法&用于抽象名词前&{surprise (惊讶) \(\rightarrow\) a surprise (一件令人惊讶的事)\\succes(成功) \(\rightarrow\) a success (一个成功的人;一件成功的事)}\\
    固定搭配&\SetCell[c=2]{l}{all of a sudden (突然地，出乎意料地), in a hurry (立即，匆忙),\\as a result of (由于), at a loss (困惑不解，茫然不知所措),\\play a key/vital/crucial role in (在……中起关键作用)}\\
  \end{talltblr}
\end{center}


\markword[\uuline]{\textbf{We}}{主语 1}
\markword[\mbox]{[currently]}{时间状语}
\markword{live}{谓语}
\markword[\mbox]{[in an era]}{时间状语}
\markword[\mbox]{{\Large\textcolor{red}{\( \lgroup \)}}}{\textcolor{red}{定从(an era)开始}}
\markword[\mbox]{in which}{引导词}
\markword[\uuline]{\textbf{technology}}{主语}
\markword{enables}{谓语}

\vspace{0.5cm}

\markword[\uwave]{information}{宾语}
\markword[\mbox]{<to reach large audiences>}{主语补足语}
\markword[\mbox]{(distributed across the globe),}{后置定语（audiences）}
\markword[\mbox]{and thus}{连词}

\vspace{0.5cm}

\markword[\uuline]{\textbf{the potential}}{主语 2}
\markword[\mbox]{(for immediate and widespread effects from misinformation)}{后置定语（potential）}

\vspace{0.5cm}

\markword[\mbox]{[now]}{时间状语}
\markword{looms}{系动词}
\markword[\uuwave]{larger}{表语}
\markword[\mbox]{[than in the past].}{比较状语}
\markword[\mbox]{{\Large\textcolor{red}{\( \rgroup \)}}}{\textcolor{red}{定从结束}}

\subsection{定冠词the}

定冠词the用法广泛，如用于序数词前（本句中的 the first 和 the third），用于专有名词前（本句中的 the Caribbean 和the Mediterranean）。现将定冠词的用法汇总如下:
\begin{center}
  \begin{talltblr}[
    label = none, % Removes the "Table X:" text
    ]{
    vlines, hlines,
    % 加入支持列表环境的测量参数 (记得在导言区加载 \UseTblrLibrary{varwidth})
    % measure=vbox, stretch=-1,
    rows={halign=c, valign=m},
    columns={halign=c, valign=m},
    column{1}={0.1\linewidth, halign=c, bg=azure9},
    column{2}={0.1\linewidth, halign=l},
    column{3}={0.3\linewidth, halign=l},
    column{4}={0.35\linewidth, halign=l},
    row{1}={font=\large\bfseries, bg=azure9, halign=c}
    }
    用法&指向&含义&举例\\
    \SetCell[r=11]{c}一般用法&\SetCell[r=2]{c}表特指&特指上文提到的人或事物&第一次提到用 a/an，以后再次提到用 the\\
        &&特指谈话双方都明确的人或物&Close the window, please. 请关上窗户。\\
        &\SetCell[r=5]{c}专用名词前&党派、民族或国家名称等专有名词&the United States 美国\\
        &&海洋、江河、湖泊、山脉、海峡、海湾等地理名词&the Yellow River 黄河\\
        &&世界上、宇宙中独一无二的事物&the earth 地球\\
        &&表示演奏乐器时，用于乐器前&play the erhu 拉二胡\\
        &&著名的历史建筑&the Great Wall 长城\\
        &\SetCell[r=2]{c}时间和方位&用于世纪或年代前&in the 1980s在20世纪 80年代\\
        &&用于方位词前面&the west 西方\\
        &\SetCell[c=2]{l}用在序数词和形容词最高级前&&the most beautiful 最漂亮的\\
        &\SetCell[c=2]{l}用于姓氏复数形式前，表示“全家人”或“夫妇俩”&&the Greens格林一家\\
    \SetCell[r=2]{c}特殊用法&\SetCell[r=2]{c}the+形容词&表示一类人（表示复数）&the poor 穷人\\
        &&表示抽象名词 (表示单数)&the right 正确的事/做法\\
    固定搭配&\SetCell[c=3]{l}{with the passage of time 随着时间的流逝\qquad{}on the grounds of 由于\\at the high expense of 以……为代价}\\
  \end{talltblr}
\end{center}


\markword[\mbox]{\Large\textcolor{orange}{\( [ \)}}{\textcolor{orange}{时壮开始}}
\markword[\mbox]{Yet from the first to the third millennium BCE, thousands of years}{}

\vspace{0.5cm}

\markword[\mbox]{\large\textcolor{olive}{\( [ \)}}{\textcolor{olive}{壮从开始}}
\markword[\mbox]{before}{引导词}
\markword[\uuline]{\textbf{these swashbucklers}}{主语}
\markword{began}{谓语}
\markword[\uwave]{spreading fear}{宾语}
\markword[\mbox]{[across the Caribbean]}{地点状语}

\vspace{0.5cm}

\markword[\mbox]{{\large\textcolor{olive}{\( ] \)}}}{\textcolor{olive}{壮从结束}}
\markword[\mbox]{{\Large\textcolor{orange}{\( ] \)}},}{\textcolor{orange}{时壮结束}}
\markword[\uuline]{\textbf{pirates}}{主语}
\markword{prowled}{谓语}
\markword[\uwave]{the Mediterranean,}{宾语}

\vspace{0.5cm}

\markword[\mbox]{[raiding merchant ships and threatening vital trade routes].}{伴随状语}

\subsection{零冠词}

有些情况下，名词前不需要添加冠词，即“零冠词”，如本句中的 at least(至少)。
\begin{center}
  \begin{talltblr}[
    label = none, % Removes the "Table X:" text
    ]{
    vlines, hlines,
    % 加入支持列表环境的测量参数 (记得在导言区加载 \UseTblrLibrary{varwidth})
    % measure=vbox, stretch=-1,
    rows={halign=c, valign=m},
    columns={halign=c, valign=m},
    column{1}={0.1\linewidth, halign=c, bg=azure9},
    column{2}={0.3\linewidth, halign=l},
    column{3}={0.45\linewidth, halign=l},
    % column{4}={0.35\linewidth, halign=l},
    row{1}={font=\large\bfseries, bg=azure9, halign=c}
    }
    用法&含义&举例\\
    \SetCell[r=5]{c}一般用法&复数名词、不可数名词前&I need help with my group assignment. 我的小组作业需要帮助。\\
        &大陆、单一的国家名称、州县、城、镇等地名前&in Australia 在澳大利亚\\
        &人名、称呼和头衔前&Sorry,sir. 抱歉，先生。\\
        &季节、月份、日期、星期几、节假日前&Spring is the best season of the year. 春天是一年中最好的季节。\\
        &球类和棋类运动、学科、语言、三餐和疾病前&I have lunch at 2 p.m. 我下午两点吃午饭。\\
    固定搭配&\SetCell[c=2]{l}{at first (起初)\qquad{}in time (及时)\qquad{}in danger (在危险中)\\lose weight (减肥)\qquad{}on time (按时)\qquad{}by car (乘坐小汽车)}\\
  \end{talltblr}
\end{center}


\markword[\uuline]{\textbf{A review}}{主语}
\markword[\mbox]{(of almost 1,800 cases)}{后置定语（review）}
\markword[\mbox]{(of entanglement in fishing nets}{}

\vspace{0.5cm}

\markword[\mbox]{and of plastic consumption among marine mammals)}{后置定语（cases）}
\markword[\mbox]{[in US waters]}{地点状语}
\markword[\mbox]{[from 2009 to 2020]}{时间状语}

\vspace{0.5cm}

\markword{found}{谓语}
\markword[\mbox]{\large\textcolor{blue}{\(\lgroup\)}}{\textcolor{blue}{宾从开始}}
\markword[\mbox]{that}{引导词}
\markword[\uuline]{\textbf{at least 700 cases}}{主语}
\markword{involved}{谓语}
\markword[\uwave]{manatees}{宾语}
\markword[\mbox]{\large\textcolor{blue}{\(\rgroup\)}}{\textcolor{blue}{宾从结束}}


%%% Local Variables:
%%% TeX-engine: xetex
%%% TeX-master: "../IELTS_300_Sentences"
%%% End:

% LocalWords:  vlines hlines halign valign xetex

% %% -*- coding: utf-8 -*-

\section{代词、介词和连词}

\subsection{人称代词}

\markword[\uuline]{\textbf{Psychologists}}{主语}
\markword{can also help}{谓语}
\markword[\uwave]{athletes}{宾语}
\markword[\mbox]{\Large\textcolor{violet}{\( < \)}}{\textcolor{violet}{宾语补足语开始}}
\markword[\mbox]{change}{引导词/主语}
\markword[\mbox]{\large\textcolor{blue}{\(\lgroup\)}}{\textcolor{blue}{宾从开始}}
\markword[\mbox]{how}{引导词}

\vspace{0.5cm}

\markword[\uuline]{\textbf{they}}{主语}
\markword{see}{谓语}
\markword[\uwave]{their physiological responses}{宾语}
\markword[\mbox]{\large\textcolor{blue}{\(\rgroup\)}}{\textcolor{blue}{宾从结束}}
\markword[\mbox]{{\Large\textcolor{violet}{\( > \)}}}{\textcolor{violet}{宾语补足语结束}}
\markword[\mbox]{<-- such as helping}{}
\markword[\mbox]{them see a higher heart rates as excitement, rather than nerves>.}{同位语}

\begin{center}
  \begin{talltblr}[
    label = none, % Removes the "Table X:" text
    ]{
    vlines, hlines,
    row{1}={font=\large\bfseries, bg=azure9},
    rows={halign=c, valign=m},
    columns={halign=c, valign=m}
    }
    \SetCell[r=2]{m}人称代词 & \SetCell[c=3]{c}单数 &          &              & \SetCell[c=3]{c}复数 &          &          \\
                             & 第一人称             & 第二人称 & 第三人称     & 第一人称             & 第二人称 & 第三人称 \\
    主格                     & I                    & you      & he, she, it  & we                   & you      & they     \\
    宾格                     & me                   & you      & him, her, it & us                   & you      & them     \\
  \end{talltblr}
\end{center}


\subsection{物主代词}

\markword[\uuline]{\textbf{They}}{主语}
\markword{operate}{谓语}
\markword[\mbox]{[with different language tools and different social tools]}{方式状语}
\markword[\mbox]{[from teachers]}{比较状语}

\vspace{0.5cm}

\markword[\mbox]{and,}{连词}
\markword[\mbox]{[haveing just learnt it themselves],}{原因状语}
\markword[\uuline]{\textbf{they}}{主语 2}
\markword{possess}{谓语 2}
\markword[\uwave]{similar cognitive structures}{宾语}

\vspace{0.5cm}

\markword[\mbox]{[to their struggling classmates].}{比较状语}

\subsection{反身代词：-self}

\subsection{指示代词this/these/that/those}

\subsection{不定代词some和any}

\subsection{不定代词each和every}

\subsection{不定代词all，both/either和neither/none}

\subsection{most和most of}

\subsection{代词it}

\subsection{关系代词who/whom/whose/that/which}

\subsection{高频介词by/in/at/on/of}

\subsection{介词短语：动词+介词 (about/at/between/for/from/of/on/to/with\(\ldots\))}

\subsection{介词短语：形容词+介词 (about/at/by/for/in/of/to/with\(\ldots\))}

\subsection{介词短语：名词+介词 (about/for/in/of/on/with\(\ldots\))}

\subsection{并列连词：并列、选择、转折、因果}

\subsection{并列连词短语：rather than}

\subsection{从属连词}

\begin{center}
  \begin{talltblr}[
    label = none, % Removes the "Table X:" text
    ]{
    vlines, hlines,
    row{1}={font=\large\bfseries, bg=azure9},
    rows={halign=c, valign=m},
    columns={halign=c, valign=m},
    column{1}={0.18\textwidth},
    column{2}={0.1\textwidth},
    column{3}={.72\textwidth, halign=l}
    }
    作用                         & \SetCell[c=2]{c}从属连词                   & \\
    引导名词性从句               & \SetCell[c=2]{l}that (无实义)；if；whether & \\
    \SetCell[r=9]{c}引导状语从句 & 时间                                       & \\
                                 & 条件                                       & \\
                                 & 目的                                       & \\
                                 & 结果                                       & \\
                                 & 原因                                       & \\
                                 & 让步                                       & \\
                                 & 方式                                       & \\
                                 & 地点                                       & \\
                                 & 比较                                       & \\
  \end{talltblr}
\end{center}


%%% Local Variables:
%%% TeX-engine: xetex
%%% TeX-master: "../IELTS_300_Sentences"
%%% End:

% LocalWords:  vlines hlines halign valign xetex LocalWords

% %% -*- coding: utf-8 -*-


\section{形容词、副词和数词}

\subsection{形容词位置：前置或后置}

\subsection{形容词顺序}

\begin{center}
  \begin{talltblr}[
    label = none, % Removes the "Table X:" text
    ]{
    vlines, hlines,
    row{1}={font=\large\bfseries, bg=azure9},
    rows={halign=c, valign=m},
    columns={halign=c, valign=m}
    }
    \SetCell[c=9]{c}形容词的排序规则                                                             \\
    数量限定词 & 观点      & 大小  & 形状  & 新旧 & 颜色 & 国籍或产地 & 材料      & 被修饰的名词 \\
    two        & beautiful & small & round & new  & blue & Chinese    & porcelain & vases        \\
    \SetCell[c=9]{c}两只小巧精美、崭新的、蓝色的中式圆瓷瓶\\
  \end{talltblr}
\end{center}

\subsection{\(-ed\)形容词和\(-ing\)形容词}

\subsection{形容词比较级和最高级的形式}

\begin{enumerate}
  \item 规则变化
  \item 不规则变化
\end{enumerate}

\subsection{形容词比较级和最高级的用法}

\subsection{形容词表达感受和观点的句式}

\subsection{时间、地点、方式及频率副词}

\subsection{程度、疑问、关系、连接及其他副词}

\subsection{副词修饰整句话}

\subsection{副词的比较级和最高级}

\subsection{比较级的强调用法}

\subsection{\(as\)比较结构}

\subsection{常用比较级的句式}

\subsection{描述相似}

\subsection{比较数量}

\subsection{“too \(\ldots\) to” 结构 }

\subsection{数词倍数表达}

\subsection{数词 \(double\)}

\subsection{描述数量：可数、不可数名词}

\subsection{\(few\)和\(little\)}




%%% Local Variables:
%%% TeX-engine: xetex
%%% TeX-master: "../IELTS_300_Sentences"
%%% End:

% LocalWords:  vlines hlines halign valign xetex LocalWords bg

% %% -*- coding: utf-8 -*-


\section{动词}

\subsection{实义动词}

实义动词又叫行为动词，具有实际动作或状态意义，可独立充当谓语。实义动词分为及物动词和不及物动词两大类：
\begin{enumerate}
  \item 及物动词:后面必须跟宾语 (动作对象)，宾语可直接跟在及物动词后，如read (阅读)、ask (询问)、give (给)、buy (买)等。如：I read a book. (read是及物动词，a book为宾语)
  \item 不及物动词:后面不能直接跟宾语，动作本身完整。如go (去)、look (看)、listen (听)、laugh (笑)、run (跑)等。如：The baby laughs. (laugh是不及物动词，后面无须接宾语)
\end{enumerate}

\subsection{连系动词(be/feel/become/keep/look/prove\(\ldots\))}

系动词本身有意义，却不能独立做句子的位于，通常需要与别的单词“连系”在一起使用，主要可以分为以下六大类：
\begin{center}
  \begin{talltblr}[
    label = none, % Removes the "Table X:" text
    ]{
    vlines, hlines,
    rows={halign=c, valign=m},
    columns={halign=c, valign=m},
    % column{1}={0.15\linewidth, halign=c, bg=azure9},
    column{2}={halign=l},
    % column{3}={0.35\linewidth, halign=l},
    % column{4}={0.35\linewidth, halign=l},
    row{1}={font=\large\bfseries, bg=azure9, halign=c}
    }
    系动词&例词\\
    状态系动词&am, is, are\\
    感官系动词&feel, smell, sound, taste\\
    变化系动词&become, grow, turn, fall, get, go, come, run\\
    持续系动词&keep, rest, remain, stay, stand\\
    表象系动词&look, seem, appear\\
    终止系动词&prove, turn out\\
  \end{talltblr}
\end{center}


\subsection{助动词(be/do/have)}

助动词协助实义动词，构成时态、语态、疑问、否定等结构：
\begin{center}
  \begin{talltblr}[
    label = none, % Removes the "Table X:" text
    ]{
    vlines, hlines,
    rows={halign=c, valign=m},
    columns={halign=c, valign=m},
    % column{1}={0.15\linewidth, halign=c, bg=azure9},
    % column{2}={0.35\linewidth, halign=l},
    column{3}={halign=l},
    % column{4}={0.35\linewidth, halign=l},
    row{1}={font=\large\bfseries, bg=azure9, halign=c}
    }
    助动词&形态&作用\\
    be&{am, is, are\\was, were,been}&{be+现在分词\(\rightarrow\)构成进行时态\\be+过去分词\(\rightarrow\)构成被动时态}\\
    do&does, did&{do not +动词原形\(\rightarrow\)构成否定式\\do+主语+动词原形\(\rightarrow\)构成一般疑问句\\do+动词原形\(\rightarrow\)用在祈使句或陈述句中，加强语气}\\
    have&has, had&{have+过去分词\(rightarrow\)构成完成时态\\have been+现在分词\(\rightarrow\)构成完成进行时态}\\
  \end{talltblr}
\end{center}

\subsection{转述动词}

转述动词用于转述他人所说的话、想法、观点、建议等内容，后面接的是that引导的宾语从句。
\begin{center}
  \begin{talltblr}[
    label = none, % Removes the "Table X:" text
    ]{
    vlines, hlines,
    rows={halign=c, valign=m},
    columns={halign=c, valign=m},
    column{1}={0.2\linewidth, halign=c, bg=azure9},
    column{2}={0.35\linewidth, halign=l},
    column{3}={0.35\linewidth, halign=l},
    % column{4}={0.35\linewidth, halign=l},
    row{1}={font=\large\bfseries, bg=azure9, halign=c}
    }
    用法&常见转述动词&举例\\
    转述动词（+that）&say, think, believe, suggest, claim, hope, imagine, doubt, expect, know, argue&She said that she would attend the meeting.\\
    转述动词+sb.+that&tell, inform, warn, remind, assure, convince, notify, persuade, advise, teach&The teacher told us that we should finish our homework.\\
    转述动词+to do&want, wish, hope, ask, demand, promise, choose, decide, agree, beg&He asked to leave early.\\
    转述动词+sb.+to do&tell, order, command, advise, encourage, invite, forbid, allow, permit, remind, urge&The doctor advised me to take more excise.\\
    转述动词+\(prep\).+\(v\).ing/\(n\).&object to, look forward to, stick to, devote oneself to, contribute to&She objects to working on weekends.\\
    转述动词+sb.+\(prep\).+\(v\).ing/\(n\).&introduce sb. to, accuse sb. of, suspect sb. of, blame sb. for, congratulate sb. on&They accused him of stealing the money.\\
    转述动词+\(v\).ing/\(n\).&enjoy, finish, mind, suggest, risk, avoid, practice, imagine, admit, deny&I enjoy reading books in my free time.\\
    转述动词+sb. (+\(n\).)&give, send, offer, lend, show, tell, teach, buy, make, find&My father bought me a new bike.\\
  \end{talltblr}
\end{center}

\subsection{第三人称单数时动词变化规则}

\begin{center}
  \begin{talltblr}[
    label = none, % Removes the "Table X:" text
    ]{
    vlines, hlines,
    rows={halign=c, valign=m},
    columns={halign=c, valign=m},
    column{1}={0.1\linewidth, halign=c, bg=azure9},
    column{2}={0.15\linewidth, halign=l},
    column{3}={0.15\linewidth, halign=l},
    column{4}={0.4\linewidth, halign=l},
    row{1}={font=\large\bfseries, bg=azure9, halign=c}
    }
    变化&词尾&公式&举例\\
    \SetCell[r=3]{c}规则变化&一般情况&词尾+s&assist \(\rightarrow\) assists; contribute\(\rightarrow\)contributes; explain\(\rightarrow\)explains\\
        &-s, -ch, -sh, -x, -o&词尾+es&research\(\rightarrow\)researches; finish\(\rightarrow\)finishes; fix\(\rightarrow\)fixes\\
        &辅音字母+y&y\(\rightarrow\)ies&reply\(\rightarrow\)replies; occupy\(\rightarrow\)occupies; hurry\(\rightarrow\)hurries\\
    不规则变化&\SetCell[c=3]{l}很多动词的第三人称单数都是不规则的，需特殊记忆，如：have\(\rightarrow\)has; do\(\rightarrow\)does&&\\
  \end{talltblr}
\end{center}

\subsection{过去式 (一般过去时) 和过去分词 (完成时态)}

\begin{center}
  \begin{talltblr}[
    label = none, % Removes the "Table X:" text
    ]{
    vlines, hlines,
    rows={halign=c, valign=m},
    columns={halign=c, valign=m},
    column{1}={0.1\linewidth, halign=c, bg=azure9},
    column{2}={0.15\linewidth, halign=l},
    column{3}={0.15\linewidth, halign=l},
    column{4}={0.4\linewidth, halign=l},
    row{1}={font=\large\bfseries, bg=azure9, halign=c}
    }
    变化&词尾&公式&举例\\
    \SetCell[r=4]{c}规则变化&一般情况&词尾+ed&{explain\(\rightarrow\)explained\(\rightarrow\)explained\\assist\(\rightarrow\)assisted\(\rightarrow\)assisted}\\
        &不发音的e&词尾+d&{create\(\rightarrow\)created\(\rightarrow\)created\\improve\(\rightarrow\)improved\(\rightarrow\)improved\\advise\(\rightarrow\)advised\(\rightarrow\)advised}\\
        &辅音字母+y&y\(\rightarrow\)ied&{occupy\(\rightarrow\)occupied\(\rightarrow\)occupied\\reply\(\rightarrow\)replied\(\rightarrow\)replied\\hurry\(\rightarrow\)hurried\(\rightarrow\)hurried}\\
        &重读闭音节&双写最后一个字母+ed&{prefer\(\rightarrow\)preferred\(\rightarrow\)preferred\\admit\(\rightarrow\)admitted\(\rightarrow\)admitted\\control\(\rightarrow\)controlled\(\rightarrow\)controlled}\\
    不规则变化&\SetCell[c=3]{l}需要特殊记忆，如：do\(\rightarrow\)did\(\rightarrow\)done; go\(\rightarrow\)went\(\rightarrow\)gone; make\(\rightarrow\)made\(\rightarrow\)made等&&\\
  \end{talltblr}
\end{center}


\subsection{现在分词 (\(-ing\))}

\begin{center}
  \begin{talltblr}[
    label = none, % Removes the "Table X:" text
    ]{
    vlines, hlines,
    rows={halign=c, valign=m},
    columns={halign=c, valign=m},
    column{1}={0.15\linewidth, halign=c, bg=azure9},
    column{2}={0.25\linewidth, halign=l},
    column{3}={0.45\linewidth, halign=l},
    % column{4}={0.35\linewidth, halign=l},
    row{1}={font=\large\bfseries, bg=azure9, halign=c}
    }
    词尾&变化公式&举例\\
    一般情况&+ing&monitor\(\rightarrow\)monitoring\\
    不发音字母e&去掉e+ing&evaluate\(\rightarrow\)evaluating; mediate\(\rightarrow\)mediating\\
    重读闭音节+辅音字母&双写辅音字母+ing&permit\(\rightarrow\)permitting; begin\(\rightarrow\)beginning; refer\(\rightarrow\)referring\\
    ie&ie\(\rightarrow\)ying&die\(\rightarrow\)dying; tie\(\rightarrow\)tying\\
  \end{talltblr}
\end{center}

\subsection{情态动词can/could的用法}

情态动词can和could通常表示能力、可能性，请求许可、给予许可等含义，具体用法如下：
\begin{description}
  \item [1. 表示能力] 两个词均表示“有能力做”，此时，could是can的过去式，比如:I can proficiently operate advanced software for data analysis.
  \item [2. 表示可能性] can和could通常用在疑问句和否定句中，表示“可能性”或“怀疑”，比如:The experiment couldn’t have failed due to such a minor error.
  \item [3. 表示许可] 当表示“请求许可”时，can和could都可使用，could并不是can的过去式，而是表达更加委婉的语气，比如: Could I get an extension for my assignment submission? 当表示“给予许可”时，一般只用can，比如:You can access the restricted area after getting the proper authorization.
\end{description}

\subsection{表示“能力”的动词}

\begin{center}
  \begin{talltblr}[
    label = none, % Removes the "Table X:" text
    ]{
    vlines, hlines,
    rows={halign=c, valign=m},
    columns={halign=c, valign=m},
    column{1}={0.1\linewidth, halign=c, bg=azure9},
    column{2}={0.25\linewidth, halign=l},
    column{3}={0.45\linewidth, halign=l},
    % column{4}={0.35\linewidth, halign=l},
    row{1}={font=\large\bfseries, bg=azure9, halign=c}
    }
    使用场景&动词&举例\\
    过去式&could, couldn't, be able to, manage to&She could speak three languages when she was a child.\\
    现在时&can, can't, be able to, manage to&I can analyze complex data sets within a short time.\\
    完成时&be able to, manage to&The team has managed to secure a large-scale investment.\\
    将来时&be able to, manage to&Our country will be able to achieve carbon neutrality in the future.\\
  \end{talltblr}
\end{center}

\subsection{情态动词表示“可能性”：must/may/might\(\ldots\)}

\begin{center}
  \begin{talltblr}[
    label = none, % Removes the "Table X:" text
    ]{
    vlines, hlines,
    rows={halign=c, valign=m},
    columns={halign=c, valign=m},
    column{1}={0.1\linewidth, halign=c, bg=azure9},
    column{2}={0.15\linewidth, halign=l},
    column{3}={0.55\linewidth, halign=l},
    % column{4}={0.35\linewidth, halign=l},
    row{1}={font=\large\bfseries, bg=azure9, halign=c}
    }
    可能性&情态动词&举例\\
    高&must&The light is on in his room. He must be at home.\\
    中&may, might, may not, might not&The conference may not attract as many participants as expected.\\
    低&can't, couldn't&They couldn't have reached the destination so far without a helicopter.\\
  \end{talltblr}
\end{center}


\subsection{情态动词呈现过去、现在和将来的“可能性”}

\begin{center}
  \begin{talltblr}[
    label = none, % Removes the "Table X:" text
    ]{
    vlines, hlines,
    rows={halign=c, valign=m},
    columns={halign=c, valign=m},
    column{1}={0.1\linewidth, halign=c, bg=azure9},
    column{2}={0.28\linewidth, halign=l},
    column{3}={0.18\linewidth, halign=l},
    column{4}={0.3\linewidth, halign=l},
    row{1}={font=\large\bfseries, bg=azure9, halign=c}
    }
    时间&构成&用法&举例\\
    \SetCell[r=2]{c}过去&may(not), might(not), could(n't), must, can't+have done&谈论过去的可能性&He may have missed the train because of the traffic jam.\\
        &may(not), might(not), could(n't), must, can't+have been+V.ing&谈论过去可能正在发生或进展中的事情&She could have been practicing the piano when you called.\\
    \SetCell[r=2]{c}现在&may(not), might(not), could(n't), must, can't+不带to的不定式&谈论当前的可能性&He might choose a different career path.\\
        &may(not), might(not), could(n't), must, can't+be+V.ing&谈论说话间可能正在发生或进展中的事情&She may be reading a book in the library.\\
    \SetCell[r=2]{c}将来&may(not), might(not), could(n't)+不带to的不定式&谈论未来的可能性或不确定性&The athlete might not participate in the next competition.\\
        &may(not), might(not), could(n't), must, can't+be+V.ing(同现在时的第二条)&谈论在未来某一时间点可能发生的事情&He may be traveling abroad next month.\\
  \end{talltblr}
\end{center}


\subsection{情态动词表示“义务性”和“必要性”}

\begin{center}
  \begin{talltblr}[
    label = none, % Removes the "Table X:" text
    ]{
    vlines, hlines,
    rows={halign=c, valign=m},
    columns={halign=c, valign=m},
    column{1}={0.15\linewidth, halign=c, bg=azure9},
    column{2}={0.45\linewidth, halign=l},
    column{3}={0.3\linewidth, halign=l},
    % column{4}={0.35\linewidth, halign=l},
    row{1}={font=\large\bfseries, bg=azure9, halign=c}
    }
    情态动词&用法&举例\\
    \SetCell[r=3]{c}must(n't)&must表达说话者主观上认为“必须”做某事，常用于标示牌、告示和印刷信息中&\SetCell[r=3]{l}We must protect endangered species to maintain ecological balance.\\
            &mustn't表示“禁止”&\\
            &must(n't)仅用于谈论现在时的义务和必要性&\\
    \SetCell[r=3]{c}have(got) to&have to强调客观条件或外部因素使得某人“不得不”做某事&\SetCell[r=3]{l}He had to stay at home to take care of his family.\\
            &got to和have to意思相近，但更口语化&\\
            &have(got) to可以谈论过去、现在和将来的义务和必要性&\\
    \SetCell[r=2]{c}need to&表示“需要”做某事，更侧重于从实际需求的角度出发&\SetCell[r=2]{l}You will need to learn some basic survival skills before going on a camping trip.\\
            &可以谈拉过去、现在和将来的义务和必要性&\\
  \end{talltblr}
\end{center}


\subsection{情态动词表示“无义务”}

\begin{center}
  \begin{talltblr}[
    label = none, % Removes the "Table X:" text
    ]{
    vlines, hlines,
    rows={halign=c, valign=m},
    columns={halign=c, valign=m},
    column{1}={0.1\linewidth, halign=c, bg=azure9},
    column{2}={0.15\linewidth, halign=l},
    column{3}={0.25\linewidth, halign=l},
    column{4}={0.3\linewidth, halign=l},
    row{1}={font=\large\bfseries, bg=azure9, halign=c}
    }
    作用&表达&用法说明&举例\\
    \SetCell[r=2]{c}谈论过去无义务&needn't have done&表示过去做了某事，但实际上没有必要做，含有“本不必做却做了”的意思&I needn't have got a taxi because the park is not far from the station.\\
        &{didn't need to\\didn't have to}&表示过去没有做某事的必要性，即“因为没有必要所以没做”&Dad picked me up from school so I didn't need to get a taxi home.\\
    谈论将来无义务&{not have to\\not need to}&&You won't have to take the final exam if you have excellent performance in daily classed.\\
  \end{talltblr}
\end{center}


\subsection{情态动词表示“提意见”}

情态动词 should(n't)和 ought (not) to 都可以用来提出意见或建议。
\begin{enumerate}
  \item should用于提出一般性的建议，意为“应该”，shouldn't是 should 的否定形式。We should protect the environment by using renewable energy.
  \item ought to表达的建议带有一定的道义上的责任或合理性，语气比should 稍强一些ought not to是ought to 的否定形式。We ought not to spread rumors about others, which can harm their reputations.
\end{enumerate}

\subsection{情态动词：will/would}

would可以看作是 will 的过去式，两个词所表达的含义类似，但用法上略有不同。
\begin{enumerate}
  \item \textbf{表达意愿}：will 表示现在的意愿，would表示过去的意愿。如:I will help you with your assignment if you need.
  \item \textbf{表示征求意见或提出请求}： 主要用于含有第二人称的疑问句中，would此时不是过去式，而是表示委语气。如:Would you like to go shopping with me?
  \item \textbf{表示习惯性和规律性}：will用于描述现在的习惯性动作或自然规律等,would用于描述过去的习惯性动作。如:When I was a child, I would go fishing with my grandfather on Sundays.
\end{enumerate}

\subsection{情态动词的替代}

\begin{enumerate}
  \item \textbf{副词替代情态动词表推测}:语气副词(如certainly、probably、possibly、perhaps、maybe 等)可替代情态动词(如must、may、might等)。比如: He is certainly at home because his car is in the garage.= He must be at home because his car is in the garage.
  \item 句型“It + be +形容词(如certain、likely、probable、possible、impossible等)+ that从句”替代情态动词表能力、概率和可能性。比如: It is impossible that he can solve such a complex problem in such a short time.=He can't solve such a complex problem in such a short time.
\end{enumerate}



%%% Local Variables:
%%% TeX-engine: xetex
%%% TeX-master: "../IELTS_300_Sentences"
%%% End:

% LocalWords:  vlines hlines halign valign xetex LocalWords sb es ies

% %% -*- coding: utf-8 -*-

\section{非谓语动词}

\subsection{不定式做主语}

\subsection{不定式作宾语}

\subsection{不定式作宾语补足语}

\subsection{不定式作定语}

\subsection{不定式作状语}

\subsection{不定式作其他句子成分}

\subsection{不带\(to\)的不定式}

\subsection{不定时的否定形式}

\subsection{疑问句 \(+\) 不定式}

\subsection{不定式常用句型}

\subsection{动名词作主语}

\subsection{代名词作宾语}

\subsection{动词后接doing和to do的区别}

\subsection{“连词+分词”结构}

\subsection{现在分词作定语}

\subsection{现在分词作状语}

\subsection{过去分词作定语}

\subsection{过去分词作状语}

\subsection{非谓语动词在句子中的作用}

\subsection{非谓语动词的被动式}





%%% Local Variables:
%%% TeX-engine: xetex
%%% TeX-master: "../IELTS_300_Sentences"
%%% End:

% LocalWords:  vlines hlines halign valign xetex LocalWords

% %% -*- coding: utf-8 -*-


\chapter{句子}

\section{基础句式及句子成分}

\subsection{主谓宾}

\begin{center}
  \begin{talltblr}[
    label = none, % Removes the "Table X:" text
    ]{
    vlines, hlines,
    % 加入支持列表环境的测量参数 (记得在导言区加载 \UseTblrLibrary{varwidth})
    % measure=vbox, stretch=-1,
    rows={halign=c, valign=m},
    columns={halign=c, valign=m},
    column{1}={0.25\linewidth, halign=c},
    column{2}={0.1\linewidth, halign=c},
    column{3}={0.15\linewidth, halign=c},
    column{4}={0.3\linewidth, halign=l},
    row{1}={font=\large\bfseries, bg=azure9, halign=c}
    }
    句子&\SetCell[c=2]{c}句子成分&&作用\\
    \SetCell[r=3]{l}Scientists conduct experiments.&主语&Scientists&句子主要说明的人或事物，也是谓语动作的执行者。\\
        &谓语&conduct&说明主语的动作，即主语“做什么”或“怎么做”。\\
        &宾语&experiments&谓语动作的承受者，即谓语动作、行为的对象。\\
  \end{talltblr}
\end{center}


\subsection{主系表}

\begin{center}
  \begin{talltblr}[
    label = none, % Removes the "Table X:" text
    ]{
    vlines, hlines,
    % 加入支持列表环境的测量参数 (记得在导言区加载 \UseTblrLibrary{varwidth})
    % measure=vbox, stretch=-1,
    rows={halign=c, valign=m},
    columns={halign=c, valign=m},
    column{1}={0.25\linewidth, halign=c},
    column{2}={0.1\linewidth, halign=c},
    column{3}={0.15\linewidth, halign=c},
    column{4}={0.3\linewidth, halign=l},
    row{1}={font=\large\bfseries, bg=azure9, halign=c}
    }
    句子&\SetCell[c=2]{c}句子成分&&作用\\
    \SetCell[r=3]{l}{The policy is a crucial measure.\\这项政策是一项关键举措。}&主语&The policy&句子主要说明的人或事物。\\
        &系动词&is&连接“主语”和“表语”，本身有一定的意义，但不能表示具体的动作。\\
        &表语&a crucial measure&描述主语的状态、性质或身份，即主语“是什么”或“怎么样”。\\
  \end{talltblr}
\end{center}


\subsection{祈使句：do/be/let\(\ldots\)}

祈使句表示请求、命令、叮嘱或者劝告等语气，通常用动词原形开头。
\begin{center}
  \begin{talltblr}[
    label = none, % Removes the "Table X:" text
    ]{
    vlines, hlines,
    % 加入支持列表环境的测量参数 (记得在导言区加载 \UseTblrLibrary{varwidth})
    % measure=vbox, stretch=-1,
    rows={halign=c, valign=m},
    columns={halign=c, valign=m},
    column{1}={0.07\linewidth, halign=c, bg=azure9},
    column{2}={0.08\linewidth, halign=l},
    column{3}={0.45\linewidth, halign=l},
    column{4}={0.25\linewidth, halign=l},
    row{1}={font=\large\bfseries, bg=azure9, halign=c}
    }
    \SetCell[c=3]{c}结构 &        &                                      & 举例            \\
    \SetCell[r=2]{c}Do型 & 肯定式 & 动词原形+（宾语）+（其他成分）       & Open the door!  \\
                         & 否定式 & Don't+动词原形+（宾语）+（其他成分） & Don't touch it. \\
    \SetCell[r=2]{c}Be型 & 肯定式 & Be+表语（名词或形容词）+（其他成分）       & Be quite!  \\
                         & 否定式 & Don't+be+表语（名词+形容词）+（其他成分） & Don't be angry with him. \\
    \SetCell[r=2]{c}Let型 & 肯定式 & Let+宾语+动词原形+（其他成分）       & Let me help you carry the box.  \\
                         & 否定式 & {Don't+let+宾语+动词原形+（其他成分）\\Let+宾语+not+动词原形+（其他成分）} & {Don't let the children play with fire.\\Let's not waste time.} \\
    \SetCell[r=2]{c}其他 & 否定式 & No+名词+V.ing, 表示“禁止做某事”       & No parking!  \\
                         & 否定式 & Never+祈使句肯定形式 & Never give up! \\
  \end{talltblr}
\end{center}


\subsection{并列句}

并列句之两个或两个以上的简单句用并列连词and、but、so、or等连接在一起后组成的句子，结构是“简单句+并列连词+简单句”，汇总如下：
\begin{center}
  \begin{talltblr}[
    label = none, % Removes the "Table X:" text
    ]{
    vlines, hlines,
    % 加入支持列表环境的测量参数 (记得在导言区加载 \UseTblrLibrary{varwidth})
    % measure=vbox, stretch=-1,
    rows={halign=c, valign=m},
    columns={halign=c, valign=m},
    % column{1}={0.15\linewidth, halign=c, bg=azure9},
    % column{2}={0.35\linewidth, halign=l},
    % column{3}={0.35\linewidth, halign=l},
    % column{4}={0.35\linewidth, halign=l},
    row{1}={font=\large\bfseries, bg=azure9, halign=c}
    }
    类型                                   & 常用并列连词                 & 用法                                           \\
    \SetCell[r=3]{c}转折、对比并列句       & but                          & “但是”                                         \\
                                           & yet                          & “但是、然而”                                   \\
                                           & while                        & “然而”                                         \\
    \SetCell[r=2]{c}因果并列句             & for                          & 表示原因，“因为”                               \\
                                           & so                           & 表示结果，“因此，所以”（so不能与becaue一起用） \\
    \SetCell[r=3]{c}选择、条件并列句       & or                           & “或者”                                         \\
                                           & neither\(\ldots\)nor\(\ldots\)         & “既不……也不……”                                 \\
                                           & either\(\ldots\)or                & “或者……或者……”                                 \\
    \SetCell[r=3]{c}并列、顺承、递进并列句 & and                          & “和；然后”                                     \\
                                           & both\(\ldots\)and                 & “不但……而且……；既……又……”                       \\
                                           & not only\(\ldots\)but (also)\(\ldots\) & “不仅……而且……”                                 \\
  \end{talltblr}
\end{center}


\subsection{主从复合句}

\begin{center}
  \begin{talltblr}[
    label = none, % Removes the "Table X:" text
    ]{
    vlines, hlines,
    % 加入支持列表环境的测量参数 (记得在导言区加载 \UseTblrLibrary{varwidth})
    % measure=vbox, stretch=-1,
    rows={halign=c, valign=m},
    columns={halign=c, valign=m},
    column{1}={0.1\linewidth, halign=c, bg=azure9},
    column{2}={0.3\linewidth, halign=l},
    column{3}={0.45\linewidth, halign=l},
    % column{4}={0.35\linewidth, halign=l},
    row{1}={font=\large\bfseries, bg=azure9, halign=c}
    }
    主从复合句&分类&举例\\
    名词性从句&{主语从句\ 宾语从句\\表语从句\ 同位语从句}&Many experts believe that AI will being profound changes to our society.\\
    定语从句&{限制性定语从句\\非限制性定语从句}&The policy which was introduced last month aims to promote economic growth.\\
    状语从句&时间、地点、原因、结果、目的、让步、条件、比较、方式\textit{状语从句}&You can find this rare plant where the climate is humid.\\
  \end{talltblr}
\end{center}

\subsection{There be句式结构}

\begin{center}
  \begin{talltblr}[
    label = none, % Removes the "Table X:" text
    ]{
    vlines, hlines,
    % 加入支持列表环境的测量参数 (记得在导言区加载 \UseTblrLibrary{varwidth})
    measure=vbox, stretch=-1,
    rows={halign=c, valign=m},
    columns={halign=c, valign=m},
    column{1}={0.2\linewidth, halign=c, bg=azure9},
    column{2}={0.65\linewidth, halign=l},
    % column{3}={0.35\linewidth, halign=l},
    % column{4}={0.35\linewidth, halign=l},
    row{1}={font=\large\bfseries, bg=azure9, halign=c}
    }
    类型&用法\\
    There be+主语+其他&{
                        \begin{enumerate}[nosep]
                          \item 通常用来表示“某处又某人或某物”，there为虚词，没有具体意义。
                          \item be为句子的谓语，be后面的名词或名词短语为句子的真正主语。
                        \end{enumerate}
                        }\\
    There be+主语+不定式&作后置定语，不定式常用主动形式表示被动意义。\\
    There+情态动词+be+主语+其他&情态动词可以是：can, may, must, should, will, 表达自己的情感。\\
    There was a time when+定语从句&意为“曾经一度，有一个时期”，用于描述个人经历。\\
  \end{talltblr}
\end{center}

\subsection{There be句式时态}

\begin{center}
  \begin{talltblr}[
    label = none, % Removes the "Table X:" text
    ]{
    vlines, hlines,
    % 加入支持列表环境的测量参数 (记得在导言区加载 \UseTblrLibrary{varwidth})
    % measure=vbox, stretch=-1,
    rows={halign=c, valign=m},
    columns={halign=c, valign=m},
    % column{1}={0.15\linewidth, halign=c, bg=azure9},
    % column{2}={0.35\linewidth, halign=l},
    % column{3}={0.35\linewidth, halign=l},
    % column{4}={0.35\linewidth, halign=l},
    row{1}={font=\large\bfseries, bg=azure9, halign=c}
    }
    时态&be动词变化形式\\
    一般现在时&There is/are…\\
    一般过去时&There was/were…\\
    一般将来时&{There will be …\\There is/are going to …}\\
    现在完成时&There has/have been …\\
    过去完成时&There had been …\\
    过去将来时&{There would be …\\There was/were going to be …}\\
  \end{talltblr}
\end{center}


\subsection{独立主格结构（特殊的状语结构）}

独立主格结构是一种特殊的状语结构，它是又名词或代词加上非谓语动词或形容词、副词、介词短语构成的。该结构有自己的“逻辑主语”，与主句的主语不同，汇总如下：
\begin{center}
  \begin{talltblr}[
    label = none, % Removes the "Table X:" text
    ]{
    vlines, hlines,
    % 加入支持列表环境的测量参数 (记得在导言区加载 \UseTblrLibrary{varwidth})
    % measure=vbox, stretch=-1,
    rows={halign=c, valign=m},
    columns={halign=c, valign=m},
    column{1}={0.25\linewidth, halign=c, bg=azure9},
    column{2}={0.55\linewidth, halign=l},
    % column{3}={0.35\linewidth, halign=l},
    % column{4}={0.35\linewidth, halign=l},
    row{1}={font=\large\bfseries, bg=azure9, halign=c}
    }
    结构&举例\\
    名词/代词+分词（现在分词/过去分词）&The task finished, they went home early.\\
    名词/动词+不定式&So many children to support, they both have to work full-time.\\
    with+名词/动词+分词/不定式&With the final exam approaching, the students are busy reviewing.\\
    There/It+being+名词/动词&It being Sunday, the library is closed.\\
  \end{talltblr}
\end{center}


\subsection{\(with\)复合结构}

\begin{enumerate}
  \item with复合结构在句中常做时间、伴随、方式、原因、条件状语等，有时还可以作后置定语，常见结构如下：
        \begin{center}
          \begin{talltblr}[
            label = none, % Removes the "Table X:" text
            ]{
            vlines, hlines,
            % 加入支持列表环境的测量参数 (记得在导言区加载 \UseTblrLibrary{varwidth})
            % measure=vbox, stretch=-1,
            rows={halign=c, valign=m},
            columns={halign=c, valign=m},
            column{1}={0.3\linewidth, halign=c, bg=azure9},
            column{2}={0.5\linewidth, halign=l},
            % column{3}={0.35\linewidth, halign=l},
            % column{4}={0.35\linewidth, halign=l},
            row{1}={font=\large\bfseries, bg=azure9, halign=c}
            }
            结构&举例\\
            with+名词/动词+介词短语&Look at the man with a book under his arm. (后置定语)\\
            with+名词/动词+形容词&The old man sat on the bench with his eyes closed. (伴随状语)\\
            with+名词/动词+副词&The plants died with the water out. (原因状语)\\
          \end{talltblr}
        \end{center}
  \item with复合结构后面还可以接动词的不同形式（to do、doing和done），用法如下：
        \begin{center}
          \begin{talltblr}[
            label = none, % Removes the "Table X:" text
            ]{
            vlines, hlines,
            % 加入支持列表环境的测量参数 (记得在导言区加载 \UseTblrLibrary{varwidth})
            % measure=vbox, stretch=-1,
            rows={halign=c, valign=m},
            columns={halign=c, valign=m},
            column{1}={0.15\linewidth, halign=c, bg=azure9},
            column{2}={0.25\linewidth, halign=l},
            column{3}={0.45\linewidth, halign=l},
            % column{4}={0.35\linewidth, halign=l},
            row{1}={font=\large\bfseries, bg=azure9, halign=c}
            }
            结构&用法&举例\\
            with+名词/动词+done&过去分词和名词/动词是被动关系，表示动作已经完成&With the problem solved, they could focus on the next task.\\
            with+名词/动词+doing&表示主动或动作正在进行&With the wind blowing strongly, the trees swayed violently.\\
            with+宾语+to do&to do和宾语是被动关系，表示尚未发生的动作&With several projects to finish this week, he has to work overtime.\\
          \end{talltblr}
        \end{center}

\end{enumerate}

\subsection{形式主语\(it\)}

It作形式主语时通常可以替代三种形式，具体如下：
\begin{enumerate}
  \item 代替不定式短语，常用于以下句型:
        \begin{itemize}
          \item “It + be +形容词+(of/for sb.) + to do sth.”,比如: It is kind of you to offer me a hand. 你能帮我一把真是太好了。
          \item “It + be +名词+to do sth,”，比如:It is your obligation to obey the law. 遵守法律是你的义务。
          \item “It takes sb. some time to do sth,”, 比如:It takes her an hour to commute to work every day. 她每天上下班通勤要花一个小时。
        \end{itemize}
  \item 代替动名词短语，常与no good、no use、useless、waste 等词连用，比如:It is useless complaining about the situation. 抱怨现状没有用。
  \item 代替主语从句，比如:It is a pity that he couldn't witness the victory. 很遗憾他没办法见证这场胜利。
\end{enumerate}

\subsection{双宾语：直接宾语和间接宾语}

宾语在句子中主要充当动作的承受者。在英语句子中，还经常会出现双宾语的情况也就是“直接宾语”和“间接宾语”，关于“双宾语”的用法汇总如下：
\begin{enumerate}
  \item 直接宾语表示动作的承受者，一般是物；间接宾语表示动作是对谁或为谁做的，一般是人。比如:I sent my friend an email.我给我的朋友发了一封电子邮件。 (其中 an emai就是直接宾语，my friend 为间接宾语)
  \item 直接宾语和间接宾语可以互换位置，比如provide sb. with sth.，互换位置后就变成provide sth. for sb.，这里互换位置需要借助介词for，还有一些谓语动词后的双宾语互换位置需要借助介词 to，如：
        \begin{itemize}
          \item book sb. sth. = book sth. for sb:为某人预定某物
          \item give sb. sth. = give sth. to sb.给某人某物
          \item lend sb. sth. = lend sth. to sb.把某物借给某人
          \item offer sb. sth. = offer sth. to sb.为某人提供某物
          \item prepare sb. sth. = prepare sth. for sb.为某人准备某物
          \item promise sb. sth = promise sth.to sb.向某人承诺某事/某物
        \end{itemize}
\end{enumerate}

\subsection{形式宾语\(it\)}

\begin{enumerate}
  \item it在作形式宾语时，通常要具备两个条件:
        \begin{enumerate}
          \item 真正的宾语是不定式、动名词或从句;
          \item 有宾语补足语。
        \end{enumerate}
  \item it作形式宾语的具体用法如下:
        \begin{enumerate}
          \item 代替不定式短语,即“make/find/think/regard\(\ldots\)+it+形容词/名词+不定式”，比如I think it essential to learn a second language.
          \item 代替动名词，比如:They regarded it no good smoking in public.
          \item 代替宾语从句，比如:We all thought it a pity that the show should have been delayed.
        \end{enumerate}
\end{enumerate}

\subsection{状语}

状语用来修饰动词、形容词、副词或整个句子，依据其表达的意义和作用的不同，主要分为下面各种状语：
\begin{center}
  \begin{longtblr}[
    label = none, % Removes the "Table X:" text
    ]{
    vlines, hlines,
    row{1}={font=\large\bfseries, bg=azure9},
    rows={halign=c, valign=m},
    columns={halign=l, valign=m}
    }
    状语     & 引导词举例                             \\
    时间状语 & when, while, as, before, after, since  \\
    地点状语   & where, in at, on                       \\
    原因状语 & because, since, as, due to             \\
    结果状语 & so/such \(\ldots\) that \(\ldots\)               \\
    目的状语 & to, so that, in order that             \\
    条件状语 & if, unless, provided that              \\
    程度状语 & 程度副词、比较级/最高级结构等          \\
    让步状语 & though, although, even though, despite \\
    方式状语 & as, as if, by                          \\
    伴随状语 & 分词短语、介词短语、形容词/名词短语    \\
    比较状语 & than, as \(\ldots\) as                      \\
  \end{longtblr}
\end{center}


\subsection{定语}

定语用来修饰名词或代词，最常见的定语类型就是“形容词定语”，如：an ancient temple。除了形容词，还有下面几类结构也可以充当定语：
\begin{center}
  \begin{talltblr}[
    label = none, % Removes the "Table X:" text
    ]{
    vlines, hlines,
    % 加入支持列表环境的测量参数 (记得在导言区加载 \UseTblrLibrary{varwidth})
    % measure=vbox, stretch=-1,
    rows={halign=c, valign=m},
    columns={halign=c, valign=m},
    column{1}={0.2\linewidth, halign=c, bg=azure9},
    column{2}={0.6\linewidth, halign=l},
    % column{3}={0.35\linewidth, halign=l},
    % column{4}={0.35\linewidth, halign=l},
    row{1}={font=\large\bfseries, bg=azure9, halign=c}
    }
    类型&举例\\
    介词短语定语&The painting in the hallway was created by a local artist.\\
    分词短语定语&The rapidly developing country faces both opportunities and challenges.\\
    不定式定语&The decision to invest in renewable energy was widely praised.\\
    定语从句&The report that highlights the importance of biodiversity was presented by the expert.\\
  \end{talltblr}
\end{center}


\subsection{同位语}

同位语常跟在名词或者代词后面，对其进行解释、说明或限定，常由名词、代词、数词或从句充当，有时也会用形容词充当，比如:
\begin{itemize}
  \item Ms. Green, an experienced journalist, wrote an in-depth report. 格林女士,一位经验丰富的记者，写了一篇深入的报道。 (名词作同位语)
  \item We ourselves are responsible for our actions. 我们自己要对我们的行为负责。 (代词作同位语)
  \item You two seem very different to me. 在我看来，你们二人有很大的不同。 (数词作同位语)
  \item Every suggestion, big or small, will be carefully considered. 每一条建议,无论重要与否，都会被认真考虑。 (形容词作同位语)
  \item The belief that hard work pays off drives him to keep going. 努力就会有回报的信念驱使他不断前进。 (从句作同位语)
\end{itemize}

\subsection{插入语}

插入语在英语中属于独立成分，起到附加解释或说明的作用，如本句中的for example。句子中常用作插入语的结构有:
\begin{enumerate}
  \item 副词或副词短语，比如 fortunately (幸运的是)、honestly (说实话)等。
  \item 形容词或形容词短语,比如hard to say (很难说)、most important of all (最重要的是)等。
  \item 介词或介词短语，比如in a word (总而言之)、in conclusion (总结)等。
  \item 不定式短语，比如to be frank (老实讲)、to tell you the truth (说实话)等。
  \item 分句，比如that is to say (也就是说)、what's more (更甚者)等。
  \item 现在分词短语，比如 generally speaking (一般来说)等。
\end{enumerate}

\subsection{直接引语和间接引语}

直接引语变间接引用时，时态变化通常遵循“时态倒退”的原则，即将直接引语中的时态向过去推移，具体变化如下：
\begin{center}
  \begin{talltblr}[
    label = none, % Removes the "Table X:" text
    ]{
    vlines, hlines,
    % 加入支持列表环境的测量参数 (记得在导言区加载 \UseTblrLibrary{varwidth})
    % measure=vbox, stretch=-1,
    rows={halign=c, valign=m},
    columns={halign=c, valign=m},
    column{1}={0.25\linewidth, halign=c, bg=azure9},
    column{2}={0.25\linewidth, halign=l},
    column{3}={0.3\linewidth, halign=l},
    % column{4}={0.35\linewidth, halign=l},
    row{1}={font=\large\bfseries, bg=azure9, halign=c}
    }
    变化&直接引语&间接引用\\
    一般现在时\(\rightarrow\)一般过去时&I am happy.&He said that he was happy.\\
    现在进行时\(\rightarrow\)过去进行时&I am reading a book.&She said that she was reading a book.\\
    一般过去时\(\rightarrow\)一般过去时或过去完成时&I saw a film yesterday.&She said that she had seen the film yesterday.\\
    过去进行时\(\rightarrow\)过去进行时或过去完成进行时&I was watching TV at that time.&He said that he was watching TV at that time.\\
    现在完成时\(\rightarrow\)过去完成时&We have finished the work.&They said that they had finished the work.\\
    过去完成时\(\rightarrow\)过去完成时&I had finished my homework before supper.&She said that she had finished her homework before supper.\\
    一般将来时\(\rightarrow\)过去将来时&I will go to the park tomorrow.&He said that he would go to the park the next day.\\
  \end{talltblr}
\end{center}
此外，直接引语变为间接引语时，还要注意人称上的变化规则：
\begin{center}
  \begin{talltblr}[
    label = none, % Removes the "Table X:" text
    ]{
    vlines, hlines,
    % 加入支持列表环境的测量参数 (记得在导言区加载 \UseTblrLibrary{varwidth})
    % measure=vbox, stretch=-1,
    rows={halign=c, valign=m},
    columns={halign=c, valign=m},
    % column{1}={0.15\linewidth, halign=c, bg=azure9},
    % column{2}={0.35\linewidth, halign=l},
    % column{3}={0.35\linewidth, halign=l},
    % column{4}={0.35\linewidth, halign=l},
    row{1}={font=\large\bfseries, bg=azure9, halign=c}
    }
    变化&特点&直接引语&间接引语\\
    \SetCell[r=3]{c}人称&\SetCell[r=3]{l}一随主，二随宾，第三人称不变化&从句主语是第一人称&从句主语和主句的人称一致\\
        &&从句主语是第二人称&从句主语和主句宾语的人称一致\\
        &&从句主语是第三人称&从句主语人称不变\\
  \end{talltblr}
\end{center}




%%% Local Variables:
%%% TeX-engine: xetex
%%% TeX-master: "../IELTS_300_Sentences"
%%% End:

% LocalWords:  vlines hlines halign valign xetex LocalWords utf sth sb

% %% -*- coding: utf-8 -*-


\section{主谓一致}

\subsection{\(and\)连接主语时的主谓一致}

\subsection{\(or\)连接主语时的主谓一致}

\subsection{\(together\) \(with\) 连接主语时的主谓一致}

\subsection{主语为单复数同形时（series/species/means\(\ldots\)）的主谓一致}

\subsection{集体名词（people/police/audience/class/family/group\(\ldots\)）的主谓一致}

\subsection{“most of \(\ldots\)”作主语时的主谓一致}

\subsection{“half of \(\ldots\)”作主语时的主谓一致}

\subsection{与number有关的主谓一致}

\subsection{分数(One third \(\ldots\))作主语时的主谓一致}

\subsection{“One of \(\ldots\)”作主语时的主谓一致}

\subsection{百分比数作主语时的主谓一致}

\subsection{不定数量的主谓一致}

\subsection{many/much的主谓一致}

\subsection{定语从句的主谓一致}

\subsection{不定代词的主谓一致}

\subsection{非谓语动词的主谓一致}

\subsection{倒装句的主谓一致}

\subsection{意义一致原则（most, the last, the rest of + 名词）}

\subsection{there be句型的主谓一致}




%%% Local Variables:
%%% TeX-engine: xetex
%%% TeX-master: "../IELTS_300_Sentences"
%%% End:

% LocalWords:  vlines hlines halign valign xetex LocalWords

% %% -*- coding: utf-8 -*-


\section{时态}

\subsection{一般现在时的结构}

\begin{center}
  \begin{talltblr}[
    label = none, % Removes the "Table X:" text
    ]{
    vlines, hlines,
    % 加入支持列表环境的测量参数 (记得在导言区加载 \UseTblrLibrary{varwidth})
    % measure=vbox, stretch=-1,
    rows={halign=c, valign=m},
    columns={halign=c, valign=m},
    column{1}={0.12\linewidth, halign=c, bg=azure9},
    column{2}={0.12\linewidth, halign=c},
    column{3}={0.25\linewidth, halign=l},
    column{4}={0.35\linewidth, halign=l},
    row{1}={font=\large\bfseries, bg=azure9, halign=c}
    }
    类型&句式&结构&举例\\
    \SetCell[r=3]{c}谓语动词是be&肯定句&主语+be（am/is/are）+其他&The new policy is of great significance for environmental protection.\\
        &否定句&主语+be（am/is/are）+not+其他&The project is not (isn't) feasible due to the lack of sufficient funds.\\
        &一般疑问句&be（am/is/are）+主语+其他?&Is the conference room available this afternoon?\\
    \SetCell[r=3]{c}谓语动词是实义动词&肯定句&主语+动词原形/三单+其他&She often participates in international academic seminars.\\
        &否定句&主语+do/does not+动词原形+其他&They don't understand the complexity of the situation.\\
        &一般疑问句&Do/Does+主语+动词原形+其他?&Does he always submit his work on time?\\
  \end{talltblr}
\end{center}


\subsection{一般现在时的标志词：always/frequently/regularly/seldom/sometimes/usually\(\ldots\)}

一般现在时用于表示经常发生的、习惯性的、现在反复出现的动作或状态，常用的时间状语有：always, frequently, regularly, seldom, sometimes, usually, every day, from time to time, now and then, once a week, every other day, as a rule（通常）等。

\subsection{一般现在时的用法}

\begin{center}
  \begin{talltblr}[
    label = none, % Removes the "Table X:" text
    ]{
    vlines, hlines,
    % 加入支持列表环境的测量参数 (记得在导言区加载 \UseTblrLibrary{varwidth})
    % measure=vbox, stretch=-1,
    rows={halign=c, valign=m},
    columns={halign=c, valign=m},
    % column{1}={0.15\linewidth, halign=c, bg=azure9},
    column{1}={0.35\linewidth, halign=l},
    column{2}={0.45\linewidth, halign=l},
    % column{4}={0.35\linewidth, halign=l},
    row{1}={font=\large\bfseries, bg=azure9, halign=c}
    }
    用法&举例\\
    描述客观事实或普遍真理&Light travels faster than sound.\\
    表示人或物的基本特征、特性及目前的状态&The newly bought sofa feels extremely comfortable.\\
    表示当前所发生的动作或存在的状态（这种状态带有一定的持续性）&My parents own a cozy coffee shop in the downtown.\\
    书报的标题、故事的叙述、小说/戏剧/电影等的情节介绍、图片的说明等&Figure 2 illustrates the complete architectural picture.\\
    表示经常发生的、习惯性的、现在反复出现的动作或状态&She usually takes a yoga class on weekends to relieve stress.\\
  \end{talltblr}
\end{center}

\subsection{一般现在时表示将来}

\begin{center}
  \begin{talltblr}[
    label = none, % Removes the "Table X:" text
    ]{
    vlines, hlines,
    % 加入支持列表环境的测量参数 (记得在导言区加载 \UseTblrLibrary{varwidth})
    % measure=vbox, stretch=-1,
    rows={halign=c, valign=m},
    columns={halign=c, valign=m},
    % column{1}={0.15\linewidth, halign=c, bg=azure9},
    column{1}={0.35\linewidth, halign=l},
    column{2}={0.45\linewidth, halign=l},
    % column{4}={0.35\linewidth, halign=l},
    row{1}={font=\large\bfseries, bg=azure9, halign=c}
    }
    用法&举例\\
    在时间状语从句和条件状语从句中，一般现在时表示将来要发生的动作&We'll start the project as soon as all the necessary materials arrive.\\
    谓语动词是come, go, arrive, leave, start, begin return, live等时，一般现在时可以表示将来发生的动作&She goes to Paris for a business trip next week.\\
    一般现在时可以和一个时间短语连用以表示已确定的、对将来的安排&The conference begins on the first day of next month.\\
  \end{talltblr}
\end{center}

\subsection{一般过去时的结构}

\begin{center}
  \begin{talltblr}[
    label = none, % Removes the "Table X:" text
    ]{
    vlines, hlines,
    % 加入支持列表环境的测量参数 (记得在导言区加载 \UseTblrLibrary{varwidth})
    % measure=vbox, stretch=-1,
    rows={halign=c, valign=m},
    columns={halign=c, valign=m},
    column{1}={0.12\linewidth, halign=c, bg=azure9},
    column{2}={0.12\linewidth, halign=l},
    column{3}={0.26\linewidth, halign=l},
    column{4}={0.35\linewidth, halign=l},
    row{1}={font=\large\bfseries, bg=azure9, halign=c}
    }
    类型&句式&结构&举例\\
    \SetCell[r=3]{c}谓语动词是be&肯定句&主语+be（was/were）+其他&She was a teacher before becoming a writer.\\
        &否定句&主语+be（was/were）+not+其他&They weren't satisfied with the results.\\
        &一般疑问句&be（Was/Were）+主语+其他&Were the experimental conditions strictly controlled?\\
    \SetCell[r=3]{c}谓语动词是实义动词&肯定句&主语+动词过去式+其他&The government implemented new policies last year.\\
        &否定句&主语+did not (didn't)+动词原形+其他&The company didn't allocate sufficient funds for product innovation.\\
        &一般疑问句&Did+主语+动词原形+其他&Did the manager approve the proposed marketing strategy?\\
  \end{talltblr}
\end{center}


\subsection{一般过去时的标志词：yesterday/just now/ago/at that time \(\ldots\)}

一般过去时表示过去某个特定时间发生的动作或存在的状态，在一般过去时的句子中，通常都要有表示过去的时间状语，如： yesterday, two days ago, last year, just now, in the old days, the other day, in 1990, at that time, once upon a time, etc.

\subsection{一般过去时的用法}

\begin{center}
  \begin{talltblr}[
    label = none, % Removes the "Table X:" text
    ]{
    vlines, hlines,
    % 加入支持列表环境的测量参数 (记得在导言区加载 \UseTblrLibrary{varwidth})
    % measure=vbox, stretch=-1,
    rows={halign=c, valign=m},
    columns={halign=c, valign=m},
    % column{1}={0.15\linewidth, halign=c, bg=azure9},
    column{1}={0.35\linewidth, halign=l},
    column{2}={0.45\linewidth, halign=l},
    % column{4}={0.35\linewidth, halign=l},
    row{1}={font=\large\bfseries, bg=azure9, halign=c}
    }
    用法&举例\\
    表示过去某个特定时间发生的动作或存在的状态&Did you attend the academic seminar last Friday?\\
    表示过去一段时间内经常发生或反复的动作&She practiced speaking English every day during her stay in the UK.\\
    讲述过去的故事或经历&Once upon a time, a beautiful girl lived in a big castle.\\
    讲述过去连续发生的动作&I woke up early, ate a breakfast, and headed to the exam center.\\
  \end{talltblr}
\end{center}


\subsection{没有时间状语的一般过去时：used to/would\(\ldots\)}

\begin{enumerate}
  \item 描述过去连续发生的动作要用过去时，这种情况下，往往没有表示过去的时间状语，而要通过上下文来理解。比如:The journalist rushed to the scene, interviewed the witnesses and wrote te report.这位记者赶到现场，采访了目击者并撰写了报道。
  \item 用used to do或 would do表示过去经常或反复发生的动作，比如:He used to take awalk in the morning.他以前常常在早晨散步。
  \item 有些情况发生的时间未清楚表明，但实际上是“刚才，刚刚”发生的，属于过去时间，应使用过去时态。常见的表述有“I didn't know/realize.”或“I forgot..”等。比如:I didn't realize the importance of this data until you mentioned it.直到你提到，我才意识到这些数据的重要性。
\end{enumerate}

\subsection{一般将来时will}

\begin{enumerate}
  \item 一般将来时will的结构
        \begin{center}
          \begin{talltblr}[
            label = none, % Removes the "Table X:" text
            ]{
            vlines, hlines,
            % 加入支持列表环境的测量参数 (记得在导言区加载 \UseTblrLibrary{varwidth})
            % measure=vbox, stretch=-1,
            rows={halign=c, valign=m},
            columns={halign=c, valign=m},
            column{1}={0.12\linewidth, halign=c, bg=azure9},
            column{2}={0.3\linewidth, halign=l},
            column{3}={0.43\linewidth, halign=l},
            % column{4}={0.35\linewidth, halign=l},
            row{1}={font=\large\bfseries, bg=azure9, halign=c}
            }
            句式&结构&举例\\
            肯定句&主语+will/shall+动词原形+其他&The government will implement more effective policies.\\
            否定句&主语+will/shall+not（won't）+动词原形+其他&The research team won't complete the experiment within the scheduled time.\\
            一般疑问句&Will/Shall+主语+动词原形+其他&Will the new teaching method improve students' academic performance?\\
          \end{talltblr}
        \end{center}
  \item 一般将来时will的用法
        \begin{center}
          \begin{talltblr}[
            label = none, % Removes the "Table X:" text
            ]{
            vlines, hlines,
            % 加入支持列表环境的测量参数 (记得在导言区加载 \UseTblrLibrary{varwidth})
            % measure=vbox, stretch=-1,
            rows={halign=c, valign=m},
            columns={halign=c, valign=m},
            % column{1}={0.15\linewidth, halign=c, bg=azure9},
            column{1}={0.45\linewidth, halign=l},
            column{2}={0.45\linewidth, halign=l},
            % column{4}={0.35\linewidth, halign=l},
            row{1}={font=\large\bfseries, bg=azure9, halign=c}
            }
            用法&举例\\
            基于个人观点或过去的经验做出预测&I think it will be extremely cold there.\\
            表达目前尚未安排、将来有可能发生的事情&She will probably travel to Europe next year.\\
            表达将来可能发生的非主观的事情或事实&The sun will rise in the east as usual tomorrow.\\
            表达说话者在说话间决定要做的事情，此时常用来提供帮助、作出承诺或给予建议&I will help you with your English if you have any problems.\\
          \end{talltblr}
        \end{center}
\end{enumerate}

\subsection{一般将来时 be going to}

\begin{enumerate}
  \item 一般将来时be going to的结构
        \begin{center}
          \begin{talltblr}[
            label = none, % Removes the "Table X:" text
            ]{
            vlines, hlines,
            % 加入支持列表环境的测量参数 (记得在导言区加载 \UseTblrLibrary{varwidth})
            % measure=vbox, stretch=-1,
            rows={halign=c, valign=m},
            columns={halign=c, valign=m},
            column{1}={0.12\linewidth, halign=c, bg=azure9},
            column{2}={0.3\linewidth, halign=l},
            column{3}={0.43\linewidth, halign=l},
            % column{4}={0.35\linewidth, halign=l},
            row{1}={font=\large\bfseries, bg=azure9, halign=c}
            }
            句式&结构&举例\\
            肯定句&主语+be (am/is/are)+going to+动词原形+其他&He is going to attend an international conference.\\
            否定句&主语+be (am/is/are)+not going to+动词原形+其他&She isn't going to waste her time on meaningless things.\\
            一般疑问句&be (Am/Is/Are)+主语+going to+动词原形+其他&Are you going to study abroad after graduation?\\
          \end{talltblr}
        \end{center}
  \item 一般将来时be going to的用法
        \begin{center}
          \begin{talltblr}[
            label = none, % Removes the "Table X:" text
            ]{
            vlines, hlines,
            % 加入支持列表环境的测量参数 (记得在导言区加载 \UseTblrLibrary{varwidth})
            % measure=vbox, stretch=-1,
            rows={halign=c, valign=m},
            columns={halign=c, valign=m},
            % column{1}={0.15\linewidth, halign=c, bg=azure9},
            column{1}={0.45\linewidth, halign=l},
            column{2}={0.45\linewidth, halign=l},
            % column{4}={0.35\linewidth, halign=l},
            row{1}={font=\large\bfseries, bg=azure9, halign=c}
            }
            用法&举例\\
            表达已经考虑过并打算实施的将来的某种行为&She is going to start a new business project this year.\\
            根据目前事实做出的预测&Look at those dark clouds. It's going to rain.\\
          \end{talltblr}
        \end{center}
\end{enumerate}

\subsection{will和be going to的用法}

\begin{enumerate}
  \item will 和 be going to 跟在 think、doubt、expect、 believe、probably、 certainly、definitely、be sure 这些单词和短语后面，表达了关于将来的某种看法。比如:
        \begin{itemize}
          \item We expect the prices will go down soon.我们期望价格很快会下降。
          \item They are definitely going to buy a new house.他们肯定会买一套新房子。
        \end{itemize}
  \item will和be going to在表达对未来的预测时，虽然都可以表示将来会发生的事情，但在使用上有一些细微差别。will 侧重于基于个人观点或过去的经验做出预测;be going to则侧重于根据目前事实做出预测。比如:
        \begin{itemize}
          \item I'm sure the new movie will be a big hit.我敢肯定这部新电影会大受欢迎。
          \item The water level in the river is rising rapidly, it is going to flood the nearby area.河水水位在迅速上升，附近地区要被淹了。
        \end{itemize}
\end{enumerate}

\subsection{表示预测将来的短语：be likely/bound/certain/estimated/predicted/set to}

除了使用 will 和 be going to，还可以用一些短语来表示将来。此类短语包括 be bound to、be certain to、be estimated to、be likely to、be predicted to、be set to等。比如:
\begin{itemize}
  \item Technological innovations are bound to reshape the future of industries.技术创新必将重塑各行业的未来。
  \item The new policy is set to bring about substantial changes in the housing market.新政策即将给房地产市场带来重大变化。
\end{itemize}

\subsection{一般将来时的标志语：tomorrow/soon/in one week/in the future/eventually\(\ldots\)}

在一般将来时的句子中，通常都会有表示将来的时间状语，如:the day after tomorrow (后天)、tonight (今天晚上)、tomorrow (明天)、this afternoon (今天下午)、next week (下周)、before long (不久后)、in ten years (十年后)、in the future (在将来)、in a little while (过一会儿)、eventually (最终)、sooner or later (迟早)等等。

\subsection{be about to表示将来时}

be about to表示即将发生的动作，结构如下：
\begin{center}
  \begin{talltblr}[
    label = none, % Removes the "Table X:" text
    ]{
    vlines, hlines,
    % 加入支持列表环境的测量参数 (记得在导言区加载 \UseTblrLibrary{varwidth})
    % measure=vbox, stretch=-1,
    rows={halign=c, valign=m},
    columns={halign=c, valign=m},
    column{1}={0.12\linewidth, halign=c, bg=azure9},
    column{2}={0.3\linewidth, halign=l},
    column{3}={0.43\linewidth, halign=l},
    % column{4}={0.35\linewidth, halign=l},
    row{1}={font=\large\bfseries, bg=azure9, halign=c}
    }
    句式&结构&举例\\
    肯定句&主语+be (am/is/are)+about to+动词原形+其他&The plane is about to leave.\\
    否定句&主语+be (am/is/are)+not about to+动词原形+其他&He isn't about to give up hie hard-earned opportunity easily.\\
    一般疑问句&be (Am/Is/Are)+主语+about to+动词原形+其他&Am I about to miss the train?\\
  \end{talltblr}
\end{center}

\subsection{现在进行时的结构}

\begin{center}
  \begin{talltblr}[
    label = none, % Removes the "Table X:" text
    ]{
    vlines, hlines,
    % 加入支持列表环境的测量参数 (记得在导言区加载 \UseTblrLibrary{varwidth})
    % measure=vbox, stretch=-1,
    rows={halign=c, valign=m},
    columns={halign=c, valign=m},
    column{1}={0.12\linewidth, halign=c, bg=azure9},
    column{2}={0.3\linewidth, halign=l},
    column{3}={0.43\linewidth, halign=l},
    % column{4}={0.35\linewidth, halign=l},
    row{1}={font=\large\bfseries, bg=azure9, halign=c}
    }
    句式&结构&举例\\
    肯定句&主语+be (am/is/are)+现在分词+其他&I am analyzing a complex data set.\\
    否定句&主语+be (am/is/are)+not+现在分词+其他&She isn't practicing the piano right now.\\
    一般疑问句&be (Am/Is/Are)+主语+现在分词+其他&Is she preparing for the upcoming exam?\\
  \end{talltblr}
\end{center}

\subsection{现在进行时的标志词：now/at the moment/listen/look\(\ldots\)}

在现在进行时中，经常用到的时间状语有now、at the moment、listen、look 以及一些暗示句子中的动作正在发生的词或句子。比如:
\begin{enumerate}
  \item 当句中出现了now、at the moment等时，表示句子要说明的是现在正在发生的事，动词应用现在进行时。
  \item 当句中出现的时间状语是 these days、this week、this month、this term 等时，如果句子所要表达的意义是在这一阶段正在发生的事，则动词应用现在进行时。
  \item 当句中出现了look、listen等暗示词时，说明后面谓语动词的动作正在发生，动词应用现在进行时。
\end{enumerate}

\subsection{现在进行时的一般用法}

\begin{center}
  \begin{talltblr}[
    label = none, % Removes the "Table X:" text
    ]{
    vlines, hlines,
    % 加入支持列表环境的测量参数 (记得在导言区加载 \UseTblrLibrary{varwidth})
    % measure=vbox, stretch=-1,
    rows={halign=c, valign=m},
    columns={halign=c, valign=m},
    % column{1}={0.15\linewidth, halign=c, bg=azure9},
    column{1}={0.45\linewidth, halign=l},
    column{2}={0.45\linewidth, halign=l},
    % column{4}={0.35\linewidth, halign=l},
    row{1}={font=\large\bfseries, bg=azure9, halign=c}
    }
    用法&举例\\
    表示说话时正在发生的动作&Stop making noise, I'm concentrating on my work.\\
    表示现阶段正在进行的动作，虽然此时此刻该动作不一定正在进行&She is writing a research paper these weeks.\\
    表示此时此刻某一个动作不断地重复&The boy is jumping with great joy at the sight of the toy car.\\
  \end{talltblr}
\end{center}

\subsection{现在进行时的特殊用法}

\begin{center}
  \begin{talltblr}[
    label = none, % Removes the "Table X:" text
    ]{
    vlines, hlines,
    % 加入支持列表环境的测量参数 (记得在导言区加载 \UseTblrLibrary{varwidth})
    % measure=vbox, stretch=-1,
    rows={halign=c, valign=m},
    columns={halign=c, valign=m},
    % column{1}={0.15\linewidth, halign=c, bg=azure9},
    column{1}={0.45\linewidth, halign=l},
    column{2}={0.45\linewidth, halign=l},
    % column{4}={0.35\linewidth, halign=l},
    row{1}={font=\large\bfseries, bg=azure9, halign=c}
    }
    用法&举例\\
    表示动作逐渐变化的过程，常用动词有get, become, come, go等&The problem is becoming more and more complicated.\\
    表示生气、赞叹或厌恶等感情色彩，常与always, constantly, continually, forever等副词连用&He is constantly making the same mistakes.\\
  \end{talltblr}
\end{center}

\subsection{不可用于现在进行时的动词：感官系动词、表示情感、心理状态的动词、拥有等}

并不是所有动词都可以用于“现在进行时”，比如感官系动词、表达情感态度或心理态度的词等通常并不出现在现在进行时的句子里，汇总如下：
\begin{center}
  \begin{talltblr}[
    label = none, % Removes the "Table X:" text
    ]{
    vlines, hlines,
    % 加入支持列表环境的测量参数 (记得在导言区加载 \UseTblrLibrary{varwidth})
    % measure=vbox, stretch=-1,
    rows={halign=c, valign=m},
    columns={halign=c, valign=m},
    column{1}={0.25\linewidth, halign=c, bg=azure9},
    column{2}={0.55\linewidth, halign=l},
    % column{3}={0.35\linewidth, halign=l},
    % column{4}={0.35\linewidth, halign=l},
    row{1}={font=\large\bfseries, bg=azure9, halign=c}
    }
    特殊多次&举例\\
    感官系动词&feel, hear, look, see, small, sound, taste\\
    表示情感态度的动词&adore, despise, dislike, enjoy, feel, hate, like, love, mind, prefer, want\\
    表示心理态度的动词&agree, assume, believe, disagree, forget, hope, know, regret, remember, suppose, think, understand\\
    表示“拥有”的动词&belong, have, own\\
    其他静态动词&appear, seem, contain, resemble (类似、相似), equal\\
  \end{talltblr}
\end{center}

\subsection{现在进行时表示将来}

现在进行时还可以表示按计划或安排即将发生的动作。
\begin{enumerate}
  \item 描述已经确定的计划安排，比如:
        \begin{itemize}
          \item She's having a meeting with her clients this afternoon.她今天下午要和客户开会。
          \item We're starting a new project next month.我们下个月要启动一个新项目。
        \end{itemize}
  \item 描述趋势预测以及流程步骤，比如:
        \begin{itemize}
          \item The raw materials are being transported to the factory next Tuesday.这些原材料将于下周二被运往工厂。
        \end{itemize}
\end{enumerate}

\subsection{将来进行时的结构}

\begin{center}
  \begin{talltblr}[
    label = none, % Removes the "Table X:" text
    ]{
    vlines, hlines,
    % 加入支持列表环境的测量参数 (记得在导言区加载 \UseTblrLibrary{varwidth})
    % measure=vbox, stretch=-1,
    rows={halign=c, valign=m},
    columns={halign=c, valign=m},
    column{1}={0.12\linewidth, halign=c, bg=azure9},
    column{2}={0.3\linewidth, halign=l},
    column{3}={0.43\linewidth, halign=l},
    % column{4}={0.35\linewidth, halign=l},
    row{1}={font=\large\bfseries, bg=azure9, halign=c}
    }
    句式&结构&举例\\
    肯定句&主语+will be+现在分词+其他&She will be attending a conference next Monday.\\
    否定句&主语+will+not+be+现在分词+其他&He won't be working on this project this weekend.\\
    一般疑问句&Will+主语+be+现在分词+其他？&Will they be travelling abroad next month?\\
  \end{talltblr}
\end{center}

\subsection{将来进行时的用法}

\begin{center}
  \begin{talltblr}[
    label = none, % Removes the "Table X:" text
    ]{
    vlines, hlines,
    % 加入支持列表环境的测量参数 (记得在导言区加载 \UseTblrLibrary{varwidth})
    % measure=vbox, stretch=-1,
    rows={halign=c, valign=m},
    columns={halign=c, valign=m},
    % column{1}={0.15\linewidth, halign=c, bg=azure9},
    column{1}={0.45\linewidth, halign=l},
    column{2}={0.45\linewidth, halign=l},
    % column{4}={0.35\linewidth, halign=l},
    row{1}={font=\large\bfseries, bg=azure9, halign=c}
    }
    用法&举例\\
    表示未来某个时刻正在进行或持续的动作&What will you be doing this time next year?\\
    表示按计划或安排将要发生的动作&They will be having a seminar next week.\\
    可用于询问别人的计划、打算（比一般将来时更委婉）&Will you be joining us for dinner?\\
    有时可以代替一般将来时使用，表示一种已经决定获肯定的动作或情况，或表示某动作将继续但还未完成&She will be giving a presentation at the meeting.\\
    与一般将来时连用时，表示将来进行时的动作晚于一般将来时所代表的动作&The construction will end next month and we will be moving into the new building.\\
  \end{talltblr}
\end{center}

\subsection{过去进行时的结构和一般用法}

过去进行时的结构是“was/were+doing”，其用法如下：
\begin{enumerate}
  \item 表示过去某段时间内持续进行的动作或者事情，常用的时间状语有 this morning、all day yesterday、last evening、at that time、when、while 等。比如:
        \begin{itemize}
          \item All day yesterday, the workers were busy renovating the old house.昨天一整天，工人们都在忙着翻新旧房子。
        \end{itemize}
  \item 表示在过去某个时间点发生的事情，比如:What were you doing at 9 p.m. last night?你昨天晚上九点在干什么?
  \item 表示过去某一动作发生时，另一动作正在进行，比如:When the phone rang, I was taking a shower.电话响的时候，我正在洗澡。
\end{enumerate}

\subsection{过去进行时的特殊用法}

过去进行时除了表示“过去正在发生”的事情，还有以下几类用法。
\begin{enumerate}
  \item 与频率副词 always、often、usually等连用，表示过去反复性或习惯性的动作，比如： You were always changing your mind.你总是改变自己的想法。
  \item 表示在过去即将发生的动作，比如:The professor asked the students if they were submitting their dissertations on time.教授问学生们是否会按时提交论文。
  \item 礼貌用法，比如:I was wondering if you could spare a few minutes to review my essay.不知道您能否抽出几分钟帮我看一下我的文章。
  \item 此外，还有一些动词不能用于过去进行时，比如agree、believe、love、meanremember、understand 等。
\end{enumerate}

\subsection{现在完成时的结构}

\begin{center}
  \begin{talltblr}[
    label = none, % Removes the "Table X:" text
    ]{
    vlines, hlines,
    % 加入支持列表环境的测量参数 (记得在导言区加载 \UseTblrLibrary{varwidth})
    % measure=vbox, stretch=-1,
    rows={halign=c, valign=m},
    columns={halign=c, valign=m},
    column{1}={0.12\linewidth, halign=c, bg=azure9},
    column{2}={0.3\linewidth, halign=l},
    column{3}={0.43\linewidth, halign=l},
    % column{4}={0.35\linewidth, halign=l},
    row{1}={font=\large\bfseries, bg=azure9, halign=c}
    }
    句式&结构&举例\\
    肯定句&主语+have/has+过去分词+其他&The researchers have made significant progress.\\
    否定句&主语+have/has+not+过去分词+其他&They haven't decided on the topic of their academic paper.\\
    一般疑问句&Have/Has+主语+过去分词+其他？&Has she participated in any international exchange program before?\\
  \end{talltblr}
\end{center}

\subsection{现在完成时的标志词：since/already/just/yet/never\(\ldots\)}

使用现在完成时的句子中常出现 already、just、yet、never、before、lately、up to now、so far、since 等标志词，其中 since 的用法较为典型，现汇总如下：
\begin{enumerate}
  \item “since +某一确定的时间点 (如具体的时间、日期)”，谓语动词常用现在完成时，且该动词应为延续性动词。比如:She has been committed to her research on AI since last January.从去年一月起，她就一直专注于人工智能方面的研究。
  \item “since+一段时间+ago”，谓语动词常用现在完成时，且该动词应为延续性动词。比如:They have kept in close touch with international scholars since three years ago.从年前开始，他们就一直与国际学者保持密切联系。
  \item “since+从句”，主句谓语动词常用现在完成时，从句谓语动词用一般过去时。比如:The city has experienced rapid development since the new mayor took office. 自从新市长上任以来，这座城市经历了快速发展。
  \item “It is/has been +一段时间+ since从句”,意为“自从……以来已经 (多久)了”。比如:It has been five months since they launched the new project.他们启动这个新项目已经五个月了。
\end{enumerate}

\subsection{现在完成时的用法}

\begin{center}
  \begin{talltblr}[
    label = none, % Removes the "Table X:" text
    ]{
    vlines, hlines,
    % 加入支持列表环境的测量参数 (记得在导言区加载 \UseTblrLibrary{varwidth})
    % measure=vbox, stretch=-1,
    rows={halign=c, valign=m},
    columns={halign=c, valign=m},
    % column{1}={0.15\linewidth, halign=c, bg=azure9},
    column{1}={0.45\linewidth, halign=l},
    column{2}={0.45\linewidth, halign=l},
    % column{4}={0.35\linewidth, halign=l},
    row{1}={font=\large\bfseries, bg=azure9, halign=c}
    }
    用法&举例\\
    表示发生在过去而对现在仍有影响的动作&I have just finished my lunch.\\
    表示从过去某事开始而延续至今的动作或状态（常与for或since等词连用）&I have lived here since 2015.\\
    在句型“This is+adj.最高级+名词+that从句”中，that从句中谓语要用现在完成时&This is the most challenging task that I have ever encountered in my research.\\
    在句型“It is first/second（或其他序数词） time that从句”中，如果强调该动作对现在的影响，that从句中谓语也要用现在完成时&It is the first time that I have been to Shanghai.\\
  \end{talltblr}
\end{center}

\subsection{时间状语在现在完成时中的位置}

\begin{center}
  \begin{talltblr}[
    label = none, % Removes the "Table X:" text
    ]{
    vlines, hlines,
    % 加入支持列表环境的测量参数 (记得在导言区加载 \UseTblrLibrary{varwidth})
    % measure=vbox, stretch=-1,
    rows={halign=c, valign=m},
    columns={halign=c, valign=m},
    column{1}={0.15\linewidth, halign=c, bg=azure9},
    column{2}={0.3\linewidth, halign=l},
    column{3}={0.4\linewidth, halign=l},
    % column{4}={0.35\linewidth, halign=l},
    row{1}={font=\large\bfseries, bg=azure9, halign=c}
    }
    时间状语位置&副词&举例\\
    在助动词和动词之间&already, just, ever, never, recently&She has just received an offer from a prestigious university.\\
    在动词之后&all my life, every day, yet, before, for ages, since 2000, since I was a child, etc.&We haven't finished the project yet.\\
  \end{talltblr}
\end{center}

\subsection{现在完成时vs一般过去时}

\begin{center}
  \begin{talltblr}[
    label = none, % Removes the "Table X:" text
    ]{
    vlines, hlines,
    % 加入支持列表环境的测量参数 (记得在导言区加载 \UseTblrLibrary{varwidth})
    % measure=vbox, stretch=-1,
    rows={halign=c, valign=m},
    columns={halign=c, valign=m},
    column{1}={0.15\linewidth, halign=c, bg=azure9},
    column{2}={0.3\linewidth, halign=l},
    column{3}={0.4\linewidth, halign=l},
    % column{4}={0.35\linewidth, halign=l},
    row{1}={font=\large\bfseries, bg=azure9, halign=c}
    }
    时态&结构&区别\\
    现在完成时&主语+have/has done+其他&表示动作已完成，强调该动作对现在造成的影响或结果\\
    一般过去时&{主语+was/were+其他\\主语+动词过去式+其他}&表示过去某个特定时间发生的动作或存在的状态，与现在没有直接联系\\
  \end{talltblr}
\end{center}

\subsection{过去完成时的结构}

\begin{center}
  \begin{talltblr}[
    label = none, % Removes the "Table X:" text
    ]{
    vlines, hlines,
    % 加入支持列表环境的测量参数 (记得在导言区加载 \UseTblrLibrary{varwidth})
    % measure=vbox, stretch=-1,
    rows={halign=c, valign=m},
    columns={halign=c, valign=m},
    column{1}={0.12\linewidth, halign=c, bg=azure9},
    column{2}={0.3\linewidth, halign=l},
    column{3}={0.43\linewidth, halign=l},
    % column{4}={0.35\linewidth, halign=l},
    row{1}={font=\large\bfseries, bg=azure9, halign=c}
    }
    句式&结构&举例\\
    肯定句&主语+had+过去分词+其他&He had completed all the assigned tasks before the deadline arrived.\\
    否定句&主语+had+not+过去分词+其他&They had not made any decision regarding the investment plan until last week.\\
    一般疑问句&Had+主语+过去分词+其他？&Had the researchers obtained sufficient data before they began the analysis?\\
  \end{talltblr}
\end{center}

\subsection{过去完成时的用法}

\begin{center}
  \begin{talltblr}[
    label = none, % Removes the "Table X:" text
    ]{
    vlines, hlines,
    % 加入支持列表环境的测量参数 (记得在导言区加载 \UseTblrLibrary{varwidth})
    % measure=vbox, stretch=-1,
    rows={halign=c, valign=m},
    columns={halign=c, valign=m},
    % column{1}={0.15\linewidth, halign=c, bg=azure9},
    column{1}={0.45\linewidth, halign=l},
    column{2}={0.45\linewidth, halign=l},
    % column{4}={0.35\linewidth, halign=l},
    row{1}={font=\large\bfseries, bg=azure9, halign=c}
    }
    用法&举例\\
    表示较早的过去，即某一时刻之前已完成的动作或状态&By the end of last month, the research team had collected all the necessary data.到上个月月底,研究团队已经收集到了所有必要的数据。\\
    表示由过去的某一时刻开始，一直延续到过去另一时间点的动作或状态 (常与for、since 等连用)&He said he had worked in that company since 2015.他说自2015年以来他就在那家公司工作。\\
    在含有定语从句的复合句中，如果叙述的是过去发生的事，先发生的动作要用过去完成时&The professor praised the students who had made remarkable progress.教授表扬了那些取得显著进步的学生。\\
    用于宾语从句:如果主句中用了know、realize、think、suppose、find、decide 等动词的一般过去式，且宾语从句中的动作先于主句的动作，通常用过去完成时&They realized they had underestimated the difficulty of the task.他们意识到他们低估了任务的难度。\\
    用来表示未曾实现的愿望或打算，此时谓语动词经常是hope、want、wish、mean、intend、attempt、think、expe等，表示“原本……、未能……”等语气&She had hoped to get a high score on the exam, but she didn't perform well.\\
  \end{talltblr}
\end{center}

\subsection{将来完成时的结构与用法}

\begin{enumerate}
  \item 将来完成时的结构
        \begin{center}
          \begin{talltblr}[
            label = none, % Removes the "Table X:" text
            ]{
            vlines, hlines,
            % 加入支持列表环境的测量参数 (记得在导言区加载 \UseTblrLibrary{varwidth})
            % measure=vbox, stretch=-1,
            rows={halign=c, valign=m},
            columns={halign=c, valign=m},
            column{1}={0.12\linewidth, halign=c, bg=azure9},
            column{2}={0.3\linewidth, halign=l},
            column{3}={0.43\linewidth, halign=l},
            % column{4}={0.35\linewidth, halign=l},
            row{1}={font=\large\bfseries, bg=azure9, halign=c}
            }
            句式&结构&举例\\
            肯定句&主语+will have+过去分词+其他&By next month, I will have finished my research.\\
            否定句&主语+will not have+过去分词+其他&We will not have finished analyzing the market data by next week.\\
            一般疑问句&Will+主语+have+过去分词+其他？&Will you have submitted your application by the deadline?\\
          \end{talltblr}
        \end{center}
  \item 将来完成时的用法
        \begin{center}
          \begin{talltblr}[
            label = none, % Removes the "Table X:" text
            ]{
            vlines, hlines,
            % 加入支持列表环境的测量参数 (记得在导言区加载 \UseTblrLibrary{varwidth})
            % measure=vbox, stretch=-1,
            rows={halign=c, valign=m},
            columns={halign=c, valign=m},
            % column{1}={0.15\linewidth, halign=c, bg=azure9},
            column{1}={0.25\linewidth, halign=l},
            column{2}={0.35\linewidth, halign=l},
            column{3}={0.3\linewidth, halign=l},
            row{1}={font=\large\bfseries, bg=azure9, halign=c}
            }
            用法&举例&标志词\\
            表示在将来某个时间之前完成的动作，并且往往对将来某一时间产生影响&We will have been married for 5 years by next month.&\SetCell[r=2]{l}将来完成时常用时间状语为“by+将来某个时间”；再使用将来完成时的句子中，by the time引导的时间状语从句中使用一般现在时\\
            表示根据某情况做出的对未来的推测，相当有must have done&They will have reached home by now considering their early start.&\\
          \end{talltblr}
        \end{center}
\end{enumerate}

\subsection{过去将来时的结构与用法}

\begin{enumerate}
  \item 过去将来时的结构
        \begin{center}
          \begin{talltblr}[
            label = none, % Removes the "Table X:" text
            ]{
            vlines, hlines,
            % 加入支持列表环境的测量参数 (记得在导言区加载 \UseTblrLibrary{varwidth})
            % measure=vbox, stretch=-1,
            rows={halign=c, valign=m},
            columns={halign=c, valign=m},
            column{1}={0.12\linewidth, halign=c, bg=azure9},
            column{2}={0.4\linewidth, halign=l},
            column{3}={0.33\linewidth, halign=l},
            % column{4}={0.35\linewidth, halign=l},
            row{1}={font=\large\bfseries, bg=azure9, halign=c}
            }
            句式&结构&举例\\
            肯定句&{主语+would+动词原形+其他\\主语+be (was/were)+going to +动词原形+其他}&She was going to visit her grandparents that weekend.\\
            否定句&{主语+would+not+动词原形+其他\\主语+be (was/were)+not+going to +动词原形+其他}&They said they would not attend the meeting.\\
            一般疑问句&{Would+主语+动词原形+其他\\be (was/were)+主语+going to +动词原形+其他}？&Was she going to study abroad?\\
          \end{talltblr}
        \end{center}
  \item 过去将来时的用法
        \begin{center}
          \begin{talltblr}[
            label = none, % Removes the "Table X:" text
            ]{
            vlines, hlines,
            % 加入支持列表环境的测量参数 (记得在导言区加载 \UseTblrLibrary{varwidth})
            % measure=vbox, stretch=-1,
            rows={halign=c, valign=m},
            columns={halign=c, valign=m},
            % column{1}={0.15\linewidth, halign=c, bg=azure9},
            column{1}={0.35\linewidth, halign=l},
            column{2}={0.5\linewidth, halign=l},
            row{1}={font=\large\bfseries, bg=azure9, halign=c}
            }
            用法&举例\\
            表示从过去某一时间看将要发生的动作或存在的状态&The report predicted that the economic situation would improve in the following year.报告预测经济形势在接下来的-年内会改善。\\
            表示过去的打算、计划或意图&They were going to have a picnic last Sunday.他们上周日原本打算去野餐。\\
          \end{talltblr}
        \end{center}
\end{enumerate}

\subsection{现在完成进行时的结构与用法}

\begin{enumerate}
  \item 现在完成进行时的结构
        \begin{center}
          \begin{talltblr}[
            label = none, % Removes the "Table X:" text
            ]{
            vlines, hlines,
            % 加入支持列表环境的测量参数 (记得在导言区加载 \UseTblrLibrary{varwidth})
            % measure=vbox, stretch=-1,
            rows={halign=c, valign=m},
            columns={halign=c, valign=m},
            column{1}={0.12\linewidth, halign=c, bg=azure9},
            column{2}={0.4\linewidth, halign=l},
            column{3}={0.33\linewidth, halign=l},
            % column{4}={0.35\linewidth, halign=l},
            row{1}={font=\large\bfseries, bg=azure9, halign=c}
            }
            句式&结构&举例\\
            肯定句&主语+have/has+been+现在分词+其他&The researcher has been conducting experiments on this new material for months.\\
            否定句&主语+have/has+not+been+现在分词+其他&He hasn't been teaching here these years.这些年他并没有一直在这儿教书。\\
            一般疑问句&have/has+主语+been+现在分词+其他？&Have you been studying for the math test today?你今天一直在准备数学测试吗?\\
          \end{talltblr}
        \end{center}
  \item 现在完成进行时的用法
        \begin{center}
          \begin{talltblr}[
            label = none, % Removes the "Table X:" text
            ]{
            vlines, hlines,
            % 加入支持列表环境的测量参数 (记得在导言区加载 \UseTblrLibrary{varwidth})
            % measure=vbox, stretch=-1,
            rows={halign=c, valign=m},
            columns={halign=c, valign=m},
            % column{1}={0.15\linewidth, halign=c, bg=azure9},
            column{1}={0.35\linewidth, halign=l},
            column{2}={0.5\linewidth, halign=l},
            row{1}={font=\large\bfseries, bg=azure9, halign=c}
            }
            用法&举例\\
            表示动作从过去某个时刻开始一直持续到现在，并有可能继续下去，强调进行的过程 (常和 for 或 since 引导的时间状语连用)&The environmentalist has been advocating for wildlife protection for a decade.这位环保主义者十年来一直在倡导野生动物保护。\\
            表示动作在一段持续的时间内多次重复&Bob has been phoning Nancy every night for the past week.这一星期以来鲍勃每天晚上都给南希打电话。\\
          \end{talltblr}
        \end{center}
\end{enumerate}

\subsection{现在完成进行时vs现在完成时}

\begin{center}
  \begin{talltblr}[
    label = none, % Removes the "Table X:" text
    ]{
    vlines, hlines,
    % 加入支持列表环境的测量参数 (记得在导言区加载 \UseTblrLibrary{varwidth})
    measure=vbox, stretch=-1,
    rows={halign=c, valign=m},
    columns={halign=c, valign=m},
    % column{1}={0.15\linewidth, halign=c, bg=azure9},
    % column{2}={0.35\linewidth, halign=l},
    column{3}={0.45\linewidth, halign=l},
    % column{4}={0.35\linewidth, halign=l},
    row{1}={font=\large\bfseries, bg=azure9, halign=c}
    }
    时态&结构&区别\\
    现在完成时&have/has done&{
                              \begin{enumerate}[nosep]
                                \item 表示动作已完成
                                \item 强调动作的结果
                              \end{enumerate}
                              }\\
    现在完成进行时&have/has been doing&{
                                        \begin{enumerate}[nosep]
                                          \item 表示动作在迄今为止的一段时间内持续进行或目前仍在进行
                                          \item 强调动作持续进行的状态
                                        \end{enumerate}
                                        }\\
  \end{talltblr}
\end{center}

\subsection{使用完成时态的句型}

\begin{enumerate}
  \item 在句型“It/This/That is+形容词最高级+that从句”中,that从句中谓语要用现在完成时比如:It is the most interesting book that I have ever read.这是我读过的最有趣的书。
  \item 在句型“It/This/That is+序数词+that 从句”中，如果强调该动作对现在的影响，that从句中谓语也要用现在完成时。比如:It is the third time that she has won the award. 这是她第三次获得这个奖项了。
\end{enumerate}

\subsection{过去完成进行时的结构}

\begin{center}
  \begin{talltblr}[
    label = none, % Removes the "Table X:" text
    ]{
    vlines, hlines,
    % 加入支持列表环境的测量参数 (记得在导言区加载 \UseTblrLibrary{varwidth})
    % measure=vbox, stretch=-1,
    rows={halign=c, valign=m},
    columns={halign=c, valign=m},
    column{1}={0.12\linewidth, halign=c, bg=azure9},
    column{2}={0.4\linewidth, halign=l},
    column{3}={0.33\linewidth, halign=l},
    % column{4}={0.35\linewidth, halign=l},
    row{1}={font=\large\bfseries, bg=azure9, halign=c}
    }
    句式&结构&举例\\
    肯定句&主语+had+been+现在分词+其他&\SetCell[r=3]{l}The architect had been designing the new museum for three years before it was approved.\\
    否定句&主语+had+not+been+现在分词+其他&\\
    一般疑问句&Had+主语+been+现在分词+其他？&\\
  \end{talltblr}
\end{center}


\subsection{过去完成进行时的用法}

\begin{center}
  \begin{talltblr}[
    label = none, % Removes the "Table X:" text
    ]{
    vlines, hlines,
    % 加入支持列表环境的测量参数 (记得在导言区加载 \UseTblrLibrary{varwidth})
    % measure=vbox, stretch=-1,
    rows={halign=c, valign=m},
    columns={halign=c, valign=m},
    % column{1}={0.15\linewidth, halign=c, bg=azure9},
    column{1}={0.35\linewidth, halign=l},
    column{2}={0.5\linewidth, halign=l},
    row{1}={font=\large\bfseries, bg=azure9, halign=c}
    }
    用法&举例\\
    表示过去某一时刻之前开始的动作或产生的状态一直延续到过去某一时刻&When the meeting started, I had been preparing the materials for an hour.会议开始时，我已经准备了一个小时的材料。\\
    用在某些从句中，表示从句的动作发在主句的动作之前而且对其有影响&I heard you had been looking for me.我听说你一直在找我。\\
  \end{talltblr}
\end{center}


%%% Local Variables:
%%% TeX-engine: xetex
%%% TeX-master: "../IELTS_300_Sentences"
%%% End:

% LocalWords:  vlines hlines halign valign xetex LocalWords

% %% -*- coding: utf-8 -*-

\section{被动语态及虚拟语气}

\subsection{一般时（现在、过去、将来）的被动语态}

\begin{center}
  \begin{talltblr}[
    label = none, % Removes the "Table X:" text
    ]{
    vlines, hlines,
    % 加入支持列表环境的测量参数 (记得在导言区加载 \UseTblrLibrary{varwidth})
    % measure=vbox, stretch=-1,
    rows={halign=c, valign=m},
    columns={halign=c, valign=m},
    % column{1}={0.15\linewidth, halign=c, bg=azure9},
    % column{2}={0.35\linewidth, halign=l},
    % column{3}={0.35\linewidth, halign=l},
    column{4}={0.25\linewidth, halign=l},
    row{1}={font=\large\bfseries, bg=azure9, halign=c}
    }
    时态&句式&构成&举例\\
    \SetCell[r=4]{c}一般现在时&肯定句&主语+am/is/are done&\SetCell[r=4]{l}The decision is made by the board.\\
        &否定句&主语+am/is/are not done&\\
        &一般疑问句&am/is/are+主语+done?&\\
        &特殊疑问句&特殊疑问词+am/is/are+主语+done?&\\
    \SetCell[r=4]{c}一般过去时&肯定句&主语+was/were done&\SetCell[r=4]{l}We were not invited to the event.\\
        &否定句&主语+was/were not done&\\
        &一般疑问句&was/were+主语+done?&\\
        &特殊疑问句&特殊疑问词+was/were+主语+done?&\\
    \SetCell[r=4]{c}一般将来时&肯定句&主语+will be done&\SetCell[r=4]{l}When will the contract be renewed?\\
        &否定句&主语+will not be done&\\
        &一般疑问句&Will+主语+be done&\\
        &特殊疑问句&特殊疑问词+will+主语+be done?&\\
  \end{talltblr}
\end{center}

\subsection{进行时（现在、过去）的被动语态}

\begin{center}
  \begin{talltblr}[
    label = none, % Removes the "Table X:" text
    ]{
    vlines, hlines,
    % 加入支持列表环境的测量参数 (记得在导言区加载 \UseTblrLibrary{varwidth})
    % measure=vbox, stretch=-1,
    rows={halign=c, valign=m},
    columns={halign=c, valign=m},
    % column{1}={0.15\linewidth, halign=c, bg=azure9},
    % column{2}={0.35\linewidth, halign=l},
    % column{3}={0.35\linewidth, halign=l},
    column{4}={0.25\linewidth, halign=l},
    row{1}={font=\large\bfseries, bg=azure9, halign=c}
    }
    时态&句式&构成&举例\\
    \SetCell[r=4]{c}现在进行时&肯定句&主语+am/is/are being done&\SetCell[r=4]{l}The data is being analyzed by the researchers.\\
        &否定句&主语+am/is/are not being done&\\
        &一般疑问句&am/is/are+主语+being done?&\\
        &特殊疑问句&特殊疑问词+am/is/are+主语+being done?&\\
    \SetCell[r=4]{c}过去进行时&肯定句&主语+was/were being done&\SetCell[r=4]{l}The room was being cleaned by him at nine yesterday morning.\\
        &否定句&主语+was/were not being done&\\
        &一般疑问句&was/were+主语+being done?&\\
        &特殊疑问句&特殊疑问词+was/were+主语+being done?&\\
  \end{talltblr}
\end{center}

\subsection{完成时（现在、过去）的被动语态}

\begin{center}
  \begin{talltblr}[
    label = none, % Removes the "Table X:" text
    ]{
    vlines, hlines,
    % 加入支持列表环境的测量参数 (记得在导言区加载 \UseTblrLibrary{varwidth})
    % measure=vbox, stretch=-1,
    rows={halign=c, valign=m},
    columns={halign=c, valign=m},
    % column{1}={0.15\linewidth, halign=c, bg=azure9},
    % column{2}={0.35\linewidth, halign=l},
    % column{3}={0.35\linewidth, halign=l},
    column{4}={0.25\linewidth, halign=l},
    row{1}={font=\large\bfseries, bg=azure9, halign=c}
    }
    时态&句式&构成&举例\\
    \SetCell[r=4]{c}现在完成时&肯定句&主语+have/has been done&\SetCell[r=4]{l}The project has not been completed yet.\\
        &否定句&主语+have/has not been done&\\
        &一般疑问句&Have/Has+主语+been done?&\\
        &特殊疑问句&特殊疑问词+have/has+主语+been done?&\\
    \SetCell[r=4]{c}过去完成时&肯定句&主语+had been done&\SetCell[r=4]{l}Had the road been blocked when you passed by?\\
        &否定句&主语+had been done&\\
        &一般疑问句&Had+主语+been done?&\\
        &特殊疑问句&特殊疑问词+had+主语+being been done?&\\
  \end{talltblr}
\end{center}

\subsection{情态动词的被动语态}

\begin{center}
  \begin{talltblr}[
    label = none, % Removes the "Table X:" text
    ]{
    vlines, hlines,
    % 加入支持列表环境的测量参数 (记得在导言区加载 \UseTblrLibrary{varwidth})
    % measure=vbox, stretch=-1,
    rows={halign=c, valign=m},
    columns={halign=c, valign=m},
    % column{1}={0.15\linewidth, halign=c, bg=azure9},
    % column{2}={0.35\linewidth, halign=l},
    % column{3}={0.35\linewidth, halign=l},
    column{4}={0.25\linewidth, halign=l},
    row{1}={font=\large\bfseries, bg=azure9, halign=c}
    }
    句式&构成&举例\\
    肯定句&主语+情态动词+be done+其他&The rules must be followed by every student.\\
    否定句&主语+情态动词+not be done+其他&The decision needn't be made immediately.\\
    一般疑问句&情态动词+主语+be done+其他?&Can the work be finished ahead of schedule?\\
    特殊疑问句&特殊疑问词+情态动词+主语+be done+其他?&Where could the lost keys be found?\\
  \end{talltblr}
\end{center}

\subsection{常用被动语态的短语}

\begin{enumerate}
  \item \textbf{常用被动语态的短语}： be appointed as被任命为…、be blamed for因为……被责备、be composed of由……组成、be equipped with装备有…、be known as 被称为…、be made of由…制成、be associated with 涉及;与……有关、be compared with与……相比、be decorated with用…装饰、be filled with 被充满…、be located in 坐落于…、be named after 以……命名
  \item \textbf{没有被动语态的短语}：break out 爆发、belong to属于……、go bad (食物等)变质、fall asleep 入睡、make progress 取得进步、keep in touch with与……保持联系、set off出发，动身、pass away 去世、take place发生、turn up 出现
\end{enumerate}

\subsection{动词主动形式表示被动含义}

\begin{center}
  \begin{talltblr}[
    label = none, % Removes the "Table X:" text
    ]{
    vlines, hlines,
    % 加入支持列表环境的测量参数 (记得在导言区加载 \UseTblrLibrary{varwidth})
    % measure=vbox, stretch=-1,
    rows={halign=c, valign=m},
    columns={halign=c, valign=m},
    column{1}={0.15\linewidth, halign=c, bg=azure9},
    column{2}={0.35\linewidth, halign=l},
    column{3}={0.35\linewidth, halign=l},
    % column{4}={0.35\linewidth, halign=l},
    row{1}={font=\large\bfseries, bg=azure9, halign=c}
    }
    情况&例词&举例\\
    部分\(vi\).&change, open, develop, improve, stop, close等&His health improves gradually.\\
    表示感官的系动词&sound, taste, feel, look等&The plan sounds feasible.\\
    由well、easily、badly等修饰的部分\(vi\).&sell, read, write, wear&The latest smartphone sells extremely well.\\
  \end{talltblr}
\end{center}

\subsection{被动语态中to的还原}

\begin{enumerate}
  \item 在不定式作宾语补足语的结构中，句子变为被动语态时，原主动语态中的宾语提前至主语位置，宾语补足语(即不定式短语)的位置不变，比如:
        \begin{itemize}
          \item The teacher encouraged the students to participate in the competition.老师鼓励学生参加比赛。 (主动)
          \item The students were encouraged to participate in the competition.学生们被鼓励参加比赛。 (被动)
        \end{itemize}
  \item 不带to的不定式作宾语补足语，在变为被动语态时，省略的to要加上，常考的动词有make、let、have、see、watch、notice、hear 等，本句中的 made to 就是此用法。再比如:
        \begin{itemize}
          \item They heard her sing in the next room.他们听见她在隔壁房间唱歌。 (主动)
          \item She was heard to sing in the next room.有人听见她在隔壁房间唱歌。 (被动)
        \end{itemize}
\end{enumerate}

\subsection{其他表示被动语态的方式}

\begin{enumerate}
  \item 被动语态除了用“be+过去分词”来表示外，还可以用以下结构来表示。
        \begin{center}
          \begin{talltblr}[
            label = none, % Removes the "Table X:" text
            ]{
            vlines, hlines,
            % 加入支持列表环境的测量参数 (记得在导言区加载 \UseTblrLibrary{varwidth})
            % measure=vbox, stretch=-1,
            rows={halign=c, valign=m},
            columns={halign=c, valign=m},
            % column{1}={0.15\linewidth, halign=c, bg=azure9},
            column{1}={0.25\linewidth, halign=l},
            column{2}={0.55\linewidth, halign=l},
            % column{4}={0.35\linewidth, halign=l},
            row{1}={font=\large\bfseries, bg=azure9, halign=c}
            }
            结构&举例\\
            get+过去分词&The window got broken during the storm.窗户在暴雨中被打破了。\\
            be+形容词+不定式&The problem is hard to solve.这个问题很难被解决。\\
            make+形容词+不定式&The teacher made the lesson easy to understand.老师把这节课讲得通俗易懂。\\
            need+动名词&The house needs painting.这栋房子需要被粉刷。\\
          \end{talltblr}
        \end{center}
  \item 使用“be+过去分词”和“get+过去分词”表示被动语态时，二者之间有一定区别。
        \begin{itemize}
          \item “be+过去分词”多用来表示一般的动作或情况;
          \item “get+过去分词”多用于表示动作的结果或动作的渐变。
        \end{itemize}
  \item “get+过去分词”结构变疑问句或否定句时，需要添加助动词do或did，比如:How did that blue door get locked?那扇蓝色的门是怎么锁上的?
\end{enumerate}

\subsection{被动语态的固定用法}

\begin{enumerate}
  \item It+被动态动词+that
        \begin{itemize}
          \item It is said that\(\ldots\) 据说…
          \item It is well known that\(\ldots\) 众所周知…
          \item It is believed that\(\ldots\) 人们相信…
          \item It is hoped that\(\ldots\) 人们希望…
          \item It is found that \(\ldots\) 人们发现…
          \item It is reported that\(\ldots\) 据报道…
        \end{itemize}
  \item 主语+被动态动词+不定式： 主语+ be asked / be believed / be considered / be estimated / be expected / be felt / be known / be meant / be reported / be said / be seen / be supposed / be thought /be understood+不定式
\end{enumerate}

\subsection{被动语态的使用场景}

\begin{center}
  \begin{talltblr}[
    label = none, % Removes the "Table X:" text
    ]{
    vlines, hlines,
    % 加入支持列表环境的测量参数 (记得在导言区加载 \UseTblrLibrary{varwidth})
    % measure=vbox, stretch=-1,
    rows={halign=c, valign=m},
    columns={halign=c, valign=m},
    % column{1}={0.15\linewidth, halign=c, bg=azure9},
    column{2}={0.55\linewidth, halign=l},
    % column{3}={0.35\linewidth, halign=l},
    % column{4}={0.35\linewidth, halign=l},
    row{1}={font=\large\bfseries, bg=azure9, halign=c}
    }
    情况&举例\\
    不知道或不必提及动作的执行者&The house was robbed last night.房子昨晚被盗了。\\
    强调动作的承受者&The statue was carved by Michelangelo.这座雕像由米开朗基罗雕刻而成。\\
    动作的执行者是无生命的事物&The power was cut off by the lightning.电力被闪电切断了。\\
    用于陈述客观事实&Coffee is grown in many countries.咖啡在许多国家都有种植。\\
  \end{talltblr}
\end{center}

\subsection{虚拟条件句——与现在情况相反}

虚拟条件句可用于表达与实际情况相反的假设，以及在这种假设下可能产生的结果，常用于提出委婉的建议、表达遗憾或描述理想化情境，其结构如下:
\begin{center}
  \begin{talltblr}[
    label = none, % Removes the "Table X:" text
    ]{
    vlines, hlines,
    % 加入支持列表环境的测量参数 (记得在导言区加载 \UseTblrLibrary{varwidth})
    % measure=vbox, stretch=-1,
    rows={halign=c, valign=m},
    columns={halign=c, valign=m},
    % column{1}={0.15\linewidth, halign=c, bg=azure9},
    column{2}={0.29\linewidth, halign=l},
    column{3}={0.35\linewidth, halign=l},
    column{1}={0.25\linewidth, halign=l},
    row{1}={font=\large\bfseries, bg=azure9, halign=c}
    }
    从句&主句&举例\\
    If+主语+动词过去式 (be 动词要用 were)&主语+would/should/could/might+动词原形&If she were in your shoes, she could handle this problem easily.如果她处在你这种情况，她能轻松处理这个问题。\\
  \end{talltblr}
\end{center}

\subsection{虚拟条件句——与将来情况相反}

虚拟条件句可用来表达与将来情况相反的假设，以及在这种假设下可能产生的结果，常用于表达建议、描述愿望以及不太可能实现的情况，其结构如下:
\begin{center}
  \begin{talltblr}[
    label = none, % Removes the "Table X:" text
    ]{
    vlines, hlines,
    % 加入支持列表环境的测量参数 (记得在导言区加载 \UseTblrLibrary{varwidth})
    % measure=vbox, stretch=-1,
    rows={halign=c, valign=m},
    columns={halign=c, valign=m},
    % column{1}={0.15\linewidth, halign=c, bg=azure9},
    column{2}={0.29\linewidth, halign=l},
    column{3}={0.35\linewidth, halign=l},
    column{1}={0.25\linewidth, halign=l},
    row{1}={font=\large\bfseries, bg=azure9, halign=c}
    }
    从句&主句&举例\\
    {if+主语+动词过去式\\if+主语+should+动词原形\\if+主语+were to+动词原形}&主语+would/should/could/might+动词原形&If everyone were to reduce their carbon footprint, the environment would improve greatly.\\
  \end{talltblr}
\end{center}

\subsection{虚拟条件句——与过去情况相反}

虚拟条件句可用来表达与过去情况相反的假设，以及在这种假设下可能产生的结果，常用于表达遗憾、后悔以及对过去情况的推测，其结构如下:
\begin{center}
  \begin{talltblr}[
    label = none, % Removes the "Table X:" text
    ]{
    vlines, hlines,
    % 加入支持列表环境的测量参数 (记得在导言区加载 \UseTblrLibrary{varwidth})
    % measure=vbox, stretch=-1,
    rows={halign=c, valign=m},
    columns={halign=c, valign=m},
    % column{1}={0.15\linewidth, halign=c, bg=azure9},
    column{2}={0.29\linewidth, halign=l},
    column{3}={0.35\linewidth, halign=l},
    column{1}={0.25\linewidth, halign=l},
    row{1}={font=\large\bfseries, bg=azure9, halign=c}
    }
    从句&主句&举例\\
    if + 主语 + had + 动词过去分词&主语 + would/should/could/might + have + 过去分词&If they had double-checked the data, they would have avoided the error.\\
  \end{talltblr}
\end{center}

\subsection{would have done}

在虚拟语气中，would have done 用于非真实条件句中，表达与过去事实相反的假设，带有惋惜、遗憾、后悔等情感。比如:If I had studied harder, I would have passed the exam.如果我当时更努力学习，我就会通过考试了。类似的结构还有:
\begin{enumerate}
  \item \textbf{should have done}: 表示过去应该做某事但实际上没有做，含有责备或遗憾的意思。
        \begin{itemize}
          \item I should have gone to the doctor earlier. Now my cold is much worse.我本应该早点去看医生的。现在我的感冒严重了。
        \end{itemize}
  \item \textbf{could have done}：表示过去本来能够做某事但实际上没有做，强调能力或可能性缺失。
        \begin{itemize}
          \item I could have gone to the concert, but I didn't have the money.我本来可以去听音乐会，但我没有钱。
        \end{itemize}
\end{enumerate}

\subsection{wish用于表示虚拟语气}

“wish+宾语从句”表示不能或不太可能实现的愿望,后面的宾语从句经常使用虚拟语气,wish 可以理解为“要是……就好了”或者“但愿……”等，从句中谓语动词的时态要根据表达的时间来确定。
\begin{enumerate}
  \item 表示过去不能实现的愿望,从句的谓语动词用过去完成时 (had+动词过去分词)，比如
        \begin{itemize}
          \item I wish I had studied harder for the exam last week.我真希望上周我为考试更努力地学习了。
        \end{itemize}
  \item 表示现在不能实现的愿望，从句的谓语动词用过去式(be动词用were)，比如:
        \begin{itemize}
          \item We wish it were sunny today.我们真希望今天是晴天。
        \end{itemize}
  \item 表示将来不能实现的愿望，用“would(could)+动词原形”，比如:
        \begin{itemize}
          \item I wish our neighbours would stop making so much noise at night. 我真希望我们的邻居晚上不要再制造这么多噪声了。
        \end{itemize}
\end{enumerate}

\subsection{错综时间条件句}

当条件状语从句表示的时间和主句表示的时间不一致时，主句与从句谓语动词的形式要根据各自发生的时间来确定。
\begin{center}
  \begin{talltblr}[
    label = none, % Removes the "Table X:" text
    ]{
    vlines, hlines,
    % 加入支持列表环境的测量参数 (记得在导言区加载 \UseTblrLibrary{varwidth})
    % measure=vbox, stretch=-1,
    rows={halign=c, valign=m},
    columns={halign=c, valign=m},
    % column{1}={0.15\linewidth, halign=c, bg=azure9},
    column{2}={0.25\linewidth, halign=l},
    column{3}={0.45\linewidth, halign=l},
    % column{4}={0.35\linewidth, halign=l},
    row{1}={font=\large\bfseries, bg=azure9, halign=c}
    }
    从句&主句&举例\\
    \SetCell[r=2]{c}与过去试试修复&与现在或现在正在发生的事不符&If I had taken that job offer last year, I would have more financial stability now.如果我去年接受了那份工作邀请，我现在的经济状况会更稳定。\\
        &与未来结果不符&If she had saved enough money last year, she would be traveling around the world next month. 如果她去年存够了钱，下个月她就会在环游世界了。\\
    \SetCell[r=2]{c}与现在事实相反&与过去事实不符&If I were more outgoing, I would have made more friends at the party last night.如果我更外向一些，昨晚在派对上我就会交到更多朋友。\\
        &与未来结果不符&If I was at home, I'd be seeing my father tomorrow because it's his birthday.如果我在家，我明天就会去看我父亲，因为是他的生日。\\
    假设的未来情况&与现在或现在正在发生的事实不符&If I wasn't meeting my teacher later, I'd at home now.\\
  \end{talltblr}
\end{center}

\subsection{表示含蓄的虚拟语气}

在虚拟语气中并不总是会出现if条件句，这时就会通过其他手段来代替条件句，如本句中的 otherwise。其主要表现形式如下:
\begin{enumerate}
  \item 用 with、without、but for、otherwise、or、under 等代替条件状语从句，比如:
        \begin{itemize}
          \item But for his timely warning, we might have been in great danger.要不是他及时警告，我们可能已经陷人极大的危险之中了。
        \end{itemize}
  \item 条件暗含在上下文中，此时主、从句可以省略其中的一个，来表示说话人的一种强烈的感情，比如:
        \begin{itemize}
          \item You could have made a more reasonable decision. 你本可以做出一个更合理的决定。
        \end{itemize}
\end{enumerate}

\subsection{宾语从句中的虚拟语气}

\begin{enumerate}
  \item “wish+宾语从句”表示不能或不太可能实现的愿望时:
        \begin{enumerate}
          \item 用一般过去时表示与现在事实相反;
          \item 用过去完成时表示与过去事实相反;
          \item 用“would/should/could/might + 动词原形”表示与将来事实相反
        \end{enumerate}
        比如：
        \begin{itemize}
          \item He wishes he had apologized to his friend earlier. 他真希望当时自己早点向朋友道歉。
        \end{itemize}
  \item 在 would rather 后的宾语从句中:
        \begin{enumerate}
          \item 用一般过去时表示与现在或将来的事实相反;
          \item 用过去完成时表示与过去事实相反
        \end{enumerate}
        比如：
        \begin{itemize}
          \item He would rather she had told him the truth earlier. 他宁愿她当时早点告诉他真相。
        \end{itemize}
  \item 在表示建议、要求、命令等的动词后的宾语从句中也经常使用虚拟语气,宾语从句中谓语动词的结构为“should+动词原形”，其中should 经常省略，比如：
        \begin{itemize}
          \item The boss demanded that the project (should) be completed within a week. 老板要求这个项目在一周内完成。
        \end{itemize}
\end{enumerate}



%%% Local Variables:
%%% TeX-engine: xetex
%%% TeX-master: "../IELTS_300_Sentences"
%%% End:

% LocalWords:  vlines hlines halign valign xetex LocalWords

% %% -*- coding: utf-8 -*-

\section{名词性从句}

\subsection{主语从句——it作形式主语（1）}

\subsection{主语从句——it作形式主语（2）}

\subsection{what引导主语从句}

\subsection{whether/if引导主语从句}

\subsection{whether/if引导宾语从句}

\subsection{宾语从句只用whether不用if}

\subsection{宾语从句连接词}

\subsection{宾语从句的that不能省略（1）}

\subsection{宾语从句的that不能省略（2）}

\subsection{宾语从句时态}

\subsection{宾语从句否定转移}

\subsection{宾语从句语序}

\subsection{表语从句连接词}

\subsection{because和why引导表语从句的区别}

\subsection{关于doubt的名词性从句}

\subsection{同位语从句修饰的抽象名词}

\subsection{同位语从句与定语从句的区别}





%%% Local Variables:
%%% TeX-engine: xetex
%%% TeX-master: "../IELTS_300_Sentences"
%%% End:

% LocalWords:  vlines hlines halign valign xetex LocalWords

% %% -*- coding: utf-8 -*-

\section{定语从句}


\subsection{定语从句引导词——关系代词（that/who/whom/which/whose/as\(\ldots\)}

\subsection{定语从句引导词that}

\subsection{定语从句引导词which和that}

\subsection{定语从句引导词whose}

\subsection{定语从句引导词who}

\subsection{介词+关系代词}

\subsection{“介词+关系代词”结构中介词的确定}

\subsection{定语从句引导词只用that而不用which的情况}

\subsection{as引导定语从句}

\subsection{定语从句引导词之关系副词（when/where/why/how\(\ldots\)）}

\subsection{定语从句引导词why}

\subsection{where表示抽象地点}

\subsection{when引导定语从句}

\subsection{非限制性定语从句的关系代词及关系副词}

\subsection{which引导的非限制性定语从句}

\subsection{两种定语（限制性和非限制性）从句的区别}

\subsection{限制性定语从句中引导词的省略}

\subsection{定语从句的特殊形式：多重和分隔式定语从句}




%%% Local Variables:
%%% TeX-engine: xetex
%%% TeX-master: "../IELTS_300_Sentences"
%%% End:

% LocalWords:  vlines hlines halign valign xetex LocalWords

% %% -*- coding: utf-8 -*-

\section{状语从句}

\subsection{when引导时间状语从句}

\subsection{表示“一……就……”意思的时间状语从句}

\subsection{as引导时间状语从句}

\subsection{until/till引导时间状语从句}

\subsection{before引导时间状语从句}

\subsection{时间状语从句的引导词：当……时，一……就……，time相关等}

\subsection{地点时间状语从句}

\subsection{because引导原因状语从句}

\subsection{since引导时间和原因状语从句}

\subsection{as引导原因状语从句}

\subsection{given that引导原因状语从句}

\subsection{其它原因状语从句常见引导词：now that/in that/seeing that\(\ldots\)}

\subsection{目的主语从句}

\subsection{in case引导目的状语从句}

\subsection{so that引导结果状语从句}

\subsection{“such \(\ldots\) that”引导结果状语从句}

\subsection{条件状语从句中“主（句）将（将来时、祈使句）从（句）现（现在时）”}

\subsection{unless引导条件状语从句}

\subsection{条件状语从句引导词：if/unless/as (so) long as/once/in case/supposing/provided\(\ldots\)}

\subsection{only if引导条件状语从句}

\subsection{if条件状语从句中的省略}

\subsection{although/though引导让步状语从句}

\subsection{even though/even if引导让步状语从句}

\subsection{“疑问词（when/where/what/which/who）\(-ever\)”引导让步状语从句}

\subsection{while的用法：引导时间、让步状语从句或引导并列句}

\subsection{常用让步状语从句引导词}

\subsection{方式状语从句：as/as if/as though}

\subsection{比较状语从句：as/than/the + 比较级}





%%% Local Variables:
%%% TeX-engine: xetex
%%% TeX-master: "../IELTS_300_Sentences"
%%% End:

% LocalWords:  vlines hlines halign valign xetex LocalWords

%% -*- coding: utf-8 -*-

\chapter{特殊语法形式}

\section{特殊用法}

\subsection{句子成分省略：主谓宾表等皆可省略}

\subsection{主语和谓语同事省略}

\subsection{不定式to的省略}

\subsection{状语从句中的省略}

\subsection{“only+状语”位于句首时的倒装形式}

\subsection{介词短语位于句首时的倒装形式}

\subsection{副词位于句首时的倒装形式}

\subsection{省略if的虚拟语气条件句中的倒装形式}

\subsection{not only的倒装形式}

\subsection{so的倒装形式}

\subsection{否地含义词（never/little/hardly/not/barely/seldom\(\ldots\)）位于句首}

\subsection{as从句的倒装形式}

\subsection{it强调句}

\subsection{强调句与其他从句的辨析}

\subsection{对“not \(\ldots\) until”结构的强调}

\subsection{what主语从句表示强调}

\subsection{其他强调方式}





%%% Local Variables:
%%% TeX-engine: xetex
%%% TeX-master: "../IELTS_300_Sentences"
%%% End:

% LocalWords:  vlines hlines halign valign xetex LocalWords

% %% -*- coding: utf-8 -*-

\section{构词法}

\subsection{表示“加强，共同”的前缀}

\begin{center}
  \begin{talltblr}[
    label = none, % Removes the "Table X:" text
    ]{
    vlines, hlines,
    row{1}={font=\large\bfseries, bg=azure9},
    rows={halign=c, valign=m},
    columns={halign=c, valign=m},
    column{2}={0.35\textwidth, halign=l},
    column{3}={0.55\textwidth, halign=l}
    }
    类别     & 前缀                                               & 举例                                                                                          \\
    表“加强” & ab-, ac-, af-, ad-, ag-, ap-, ar-, as-, at-        & abound, acclaim, affirm, adhere, aggravate, appreciate, arouse, assemble, attract             \\
    表“共同” & co-, col-, com-, con-, cor-, sym-, syn-, syl-, sy- & cooperate, collaborate, combine, connect, correspond, sympathy, synchronize, syllable, system \\
  \end{talltblr}
\end{center}

\subsection{表示否定意义的前缀}

\subsection{表示“空间、方位、场所”的前缀}

\subsection{与“数量、数字”有关的前缀}

\subsection{表示“学术、主义”的后缀}

\subsection{表示“性质、行为、状态”的后缀}

\subsection{表“人或物”的后缀}

\subsection{表示“做，使……”的后缀}

\subsection{表示“写、画”的词根}

\subsection{与“自然、科学”有关的词根}

\subsection{与“时间、空间、顺序”有关的词根}

\subsection{与“想法、观点”有关的词根}

\subsection{合成动词}

\begin{center}
  \begin{talltblr}[
    label = none, % Removes the "Table X:" text
    ]{
    vlines, hlines,
    row{1}={font=\large\bfseries, bg=azure9},
    rows={halign=c, valign=m},
    columns={halign=c, valign=m}
    }
    结构        & 举例                 \\
    名词+动词   & baby-sit, proofread  \\
    形容词+动词 & whitewash, safeguard \\
    副词+动词   & overshadow, undergo  \\
  \end{talltblr}
\end{center}

\subsection{合成名词}

\subsection{合成形容词}



%%% Local Variables:
%%% TeX-engine: xetex
%%% TeX-master: "../IELTS_300_Sentences"
%%% End:

% LocalWords:  vlines hlines halign valign xetex LocalWords

% %% -*- coding: utf-8 -*-

\section{复杂句}





%%% Local Variables:
%%% TeX-engine: xetex
%%% TeX-master: "../IELTS_300_Sentences"
%%% End:

% LocalWords:  vlines hlines halign valign xetex LocalWords



\end{document}

% Local Variables:
% TeX-engine: xetex
% TeX-master: t
% End:

% LocalWords:  AutoFallBack AutoFakeBold PunctStyle kaiming CheckFullRight GBK
% LocalWords:  CheckSingle AllowBreakBetweenPuncts CJKmath BoldFont FZYouHei li
% LocalWords:  NSimSun IosevkaTerm NFM FZZJ HMLSJF FZNewKai kai SimSun ExtB
% LocalWords:  CircNum xetex LocalWords
